%% 3-3 関数は「入れると出てくる」ルール（第12回編集委員会：組 V3-2 が書いた）  % noruby
%% 流れ：きまり当てゲーム → ① 1つ決まると1つ決まる（関数）→ ② 表・グラフ・式（華氏と摂氏）→ ③ f(x) の箱  % noruby
%%       → ④ かっこごと入れる f(a+h) → ⑤ 箱をつなぐ（合成関数）→ ⑥ いろいろな関数（負の比例・変域・階段・比例式）  % noruby
%% 第5巻 5-1・5-3・5-6 の思い出そうが指す文（f(3)、f(a+h)、u=3x+1 と y=u^2）は、この形のまま変えないこと。  % noruby
\makeatletter
% 華氏の温度  \bkthrdF{数}  → 86°F  % noruby
\newcommand\bkthrdF[1]{\mbox{\qty{#1}{\degree}F}}
% 入れると出てくる箱（機械の形）  \bkthrdbox{(x,y)}{名前}{中に書く文字}{はば}  名前は箱の上の札  % noruby
\newcommand\bkthrdbox[4]{\begin{scope}[shift={#1}]
  \fill[bkTheme!14, draw=bkDeep, line width=1pt, rounded corners=2.5mm] ({-(#4)/2},-9) rectangle ({(#4)/2},9);
  \fill[bkDeep] ({-(#4)/2-1.4},-3) rectangle ({-(#4)/2},3);
  \fill[bkDeep] ({(#4)/2},-3) rectangle ({(#4)/2+1.4},3);
  \node[font=\sffamily\small, text=bkInk, align=center] at (0,-1.6) {#3};
  \node[fill=bkTheme, draw=bkDeep, line width=0.6pt, text=bkOn, font=\sffamily\bfseries\small, rounded corners=1mm, inner xsep=2mm, inner ysep=0.8mm] at (0,9) {#2};
\end{scope}}
% 式を、まん中に1行で置く  \bkthrdline{…}  % noruby
\newcommand\bkthrdline[1]{\par\addvspace{3pt}{\centering #1\par}\addvspace{3pt}}
% 説明の段の見出し（丸数字＋太字の1行）  \bkthrdstep{番号}{見出し}  % noruby
\newcommand\bkthrdstep[2]{\par\addvspace{6pt}\Needspace{4\baselineskip}\noindent
  {\bkcircled[bkOrange]{#1}\enspace\sffamily\bfseries\color{bkOrange!75!black}#2}\par\nobreak\vspace{2pt}}
% \R{小}{ちい}さなカード（はてな？の表・グラフ・式）  \bkthrdcard{(x,y)}{はば}{高さ}{色}  % noruby
\newcommand\bkthrdcard[4]{\fill[white, draw=#4, line width=0.9pt, rounded corners=2mm] #1 rectangle ++(#2,#3);}
% 例題の「しき」の行（折り返した行を字下げする）  % noruby
%   \bkthrdS … しきの1行め（折り返しは2.6字下げ）  \bkthrdSi … しきの1行めが（1）で始まるとき（折り返しは5.1字下げ）  % noruby
%   \bkthrdI … （2）などで始まる次の行  \bkthrdJ … 番号のない続きの行（5.1字下げ）  \bkthrdC … 番号のない続きの行（2.6字下げ）  % noruby
\newcommand\bkthrdS{\shiki\hangindent=2.6\zw\hangafter=1\relax}
\newcommand\bkthrdSi{\shiki\hangindent=5.1\zw\hangafter=1\relax}
\newcommand\bkthrdI{\par\noindent\hangindent=5.1\zw\hangafter=1\relax\hspace*{2.6\zw}}
\newcommand\bkthrdJ{\par\noindent\hangindent=5.1\zw\hangafter=0\relax}
\newcommand\bkthrdC{\par\noindent\hangindent=2.6\zw\hangafter=0\relax}
\newcommand\bkthrdK{\kotae\hangindent=2.6\zw\hangafter=1\relax}
% プログラムの枠と実行結果の枠（2-19 の \bkiitcode・\bkiitout と\R{同}{おな}じ形）  % noruby
{\catcode`\ =\active\gdef\bkthrd@obey{\catcode`\ =\active\def {\leavevmode\nobreak\ }\obeylines\parindent=0pt\parskip=0pt}}
\newtcolorbox{bkthrdcode}[1][Python]{enhanced, colback=white, colframe=bkNavy!60, boxrule=0.7pt, arc=1.5mm,
  left=3mm, right=3mm, top=1.2mm, bottom=1.2mm, before skip=3pt, after skip=3pt,
  fontupper=\ttfamily\gtfamily\small, before upper={\bkthrd@obey\baselineskip=14.5pt},
  attach boxed title to top right={xshift=-3mm, yshift*=-\tcboxedtitleheight/2},
  boxed title style={colback=bkNavy!60, colframe=bkNavy!60, arc=1mm, boxrule=0pt, left=1.5mm, right=1.5mm, top=0.3mm, bottom=0.3mm},
  fonttitle=\sffamily\bfseries\scriptsize, coltitle=white, title={#1}}
\newtcolorbox{bkthrdout}[1][\R{実行}{じっ|こう}\R{結果}{けっ|か}]{enhanced, colback=bkInk!85, colframe=bkInk!85, boxrule=0.7pt, arc=1.5mm,
  left=3mm, right=3mm, top=1.2mm, bottom=1.2mm, before skip=3pt, after skip=3pt, coltext=white,
  fontupper=\ttfamily\gtfamily\small, before upper={\bkthrd@obey\baselineskip=14.5pt},
  attach boxed title to top right={xshift=-3mm, yshift*=-\tcboxedtitleheight/2},
  boxed title style={colback=bkGray, colframe=bkGray, arc=1mm, boxrule=0pt, left=1.5mm, right=1.5mm, top=0.3mm, bottom=0.3mm},
  fonttitle=\sffamily\bfseries\scriptsize, coltitle=white, title={#1}}
\makeatother

\bkchapter{3}{3}{\R{関数}{かん|すう}は\\「\R{入}{い}れると\R{出}{で}てくる」ルール}{\bklevelrange{1}{2}}{\R{関数}{かん|すう}は「\R{入}{い}れると\R{出}{で}てくる」ルール}
\bkmichishirube{2-11,2-13,2-15,3-1,3-2}{3-4,3-5,5-1,5-3,5-6,9-8}

\begin{goal}
\begin{itemize}[label={\color{\bkTT}●}, leftmargin=1.6\zw, itemsep=0pt, topsep=1pt]
\item \bklevel{1}\enspace 1つ\R{決}{き}まるともう1つが\R{決}{き}まる2つの\R{量}{りょう}（\R{関数}{かん|すう}）を\R{見分}{み|わ}け、\R{表}{ひょう}・グラフ・\R{式}{しき}で\R{表}{あらわ}せる。
\item \bklevel{2}\enspace $f(x)$ の\R{書}{か}き\R{方}{かた}を\R{読}{よ}み、$f(3)$ や $f(a+h)$ を\R{計算}{けい|さん}できる。
\item \bklevel{2}\enspace 2つの\R{関数}{かん|すう}をつないだ\R{合成}{ごう|せい}\R{関数}{かん|すう} $f(g(x))$ を\R{作}{つく}れる。
\end{itemize}
\end{goal}

\begin{hatena}
アメリカの\R{天気予報}{てん|き|よ|ほう}で「あしたの\R{最高気温}{さい|こう|き|おん}は86\R{度}{ど}」。\R{暑}{あつ}いのかな？　\R{寒}{さむ}いのかな？　アメリカでは、\R{温度}{おん|ど}をふつう\textbf{\R{華氏}{か|し}}（°F）という\R{目}{め}もりで\R{表}{あらわ}します。\R{日本}{に|ほん}で\R{使}{つか}う\textbf{\R{摂氏}{せっ|し}}（℃）とは、\R{目}{め}もりがちがうのです。3\R{人}{にん}が、\R{摂氏}{せっ|し}と\R{華氏}{か|し}の\R{関係}{かん|けい}を、それぞれのやり\R{方}{かた}で\R{調}{しら}べてきました。3\R{人}{にん}は、\R{同}{おな}じことを\R{言}{い}っているのかな？　\bkthrdF{86}は、\R{摂氏}{せっ|し}で\R{何}{なん}℃かな？
\begin{bkfigure}
\begin{tikzpicture}[x=1mm, y=1mm]
  % ソラさんの表
  \bkthrdcard{(0,0)}{40}{28}{bkBlue}
  \node[font=\sffamily\bfseries\footnotesize, text=bkBlue!80!black] at (20,23.5) {ソラさんの\R{表}{ひょう}};
  \draw[bkLine, line width=0.5pt] (3,12.5) -- (37,12.5) (11,19.5) -- (11,5);
  \node[font=\sffamily\footnotesize] at (7,16) {℃};
  \node[font=\sffamily\footnotesize] at (7,8.5) {°F};
  \foreach \c/\f [count=\i] in {0/32,10/50,20/68,100/212} {
    \node[font=\footnotesize] at (8.5+\i*6.4,16) {\c};
    \node[font=\footnotesize] at (8.5+\i*6.4,8.5) {\f};}
  % ミナさんのグラフ
  \bkthrdcard{(42,0)}{38}{28}{bkGreen!80!black}
  \node[font=\sffamily\bfseries\footnotesize, text=bkGreen!50!black] at (61,23.5) {ミナさんのグラフ};
  \draw[-{Stealth[length=3pt]}, bkInk, line width=0.6pt] (50,5) -- (76,5);
  \draw[-{Stealth[length=3pt]}, bkInk, line width=0.6pt] (50,5) -- (50,19.5);
  \draw[bkBlue, line width=1.2pt] (50,8) -- (72,17.5);
  \node[font=\sffamily\scriptsize, anchor=north] at (74.5,4.3) {℃};
  \node[font=\sffamily\scriptsize, anchor=east] at (49.3,18) {°F};
  % デルタさんの式
  \bkthrdcard{(82,0)}{40}{28}{bkOrange}
  \node[font=\sffamily\bfseries\footnotesize, text=bkOrange!80!black] at (102,23.5) {デルタさんの\R{式}{しき}};
  \node[font=\normalsize] at (102,14) {$F=1.8C+32$};
  \node[font=\sffamily\scriptsize, text=bkGray, align=center] at (102,6) {$C$ が\R{摂氏}{せっ|し}、$F$ が\R{華氏}{か|し}};
\end{tikzpicture}
\end{bkfigure}
\end{hatena}

\clearpage
%% ---------------- やってみよう・思い出そう ----------------
\begin{yattemiyou}[きまり\R{当}{あ}てゲーム]
\R{2人}{ふたり}1\R{組}{くみ}で\R{遊}{あそ}びます。\R{1人}{ひとり}が「\R{出}{だ}す\R{人}{ひと}」、もう\R{1人}{ひとり}が「\R{当}{あ}てる\R{人}{ひと}」です。
\begin{enumerate}[label={\bkcircled[bkGreen]{\arabic*}}, leftmargin=2.4\zw, itemsep=1pt, topsep=2pt]
\item \R{出}{だ}す\R{人}{ひと}は、「2\R{倍}{ばい}して3たす」のような\R{計算}{けい|さん}のきまりを、\R{紙}{かみ}にこっそり\R{書}{か}いておく。
\item \R{当}{あ}てる\R{人}{ひと}は、\R{好}{す}きな\R{数}{かず}を1つ\R{言}{い}う。\R{出}{だ}す\R{人}{ひと}は、きまりのとおりに\R{計算}{けい|さん}した\R{数}{かず}を\R{答}{こた}える。
\item \R{当}{あ}てる\R{人}{ひと}は、\R{言}{い}った\R{数}{かず}と\R{答}{こた}えの\R{数}{かず}を、\R{表}{ひょう}に\R{書}{か}いていく。
\item きまりがわかったら、\R{当}{あ}てる\R{人}{ひと}は「\R{入}{い}れた\R{数}{かず}を $x$、\R{出}{で}てきた\R{数}{かず}を $y$ とすると、$y=\cdots$」と\R{式}{しき}で\R{答}{こた}えよう。
\end{enumerate}
\par\vspace{3pt}
ソラさんが\R{当}{あ}てる\R{人}{ひと}になったときの\R{表}{ひょう}です。どんなきまりかな？
\begin{bkcenter}
\begin{bktable}{colspec={Q[l,m,wd=32mm] *{5}{Q[c,m,wd=11mm]}}, rowsep=1.5pt}
\R{言}{い}った\R{数}{かず} $x$ & 1 & 2 & 5 & 10 & 0 \\
\R{答}{こた}えの\R{数}{かず} $y$ & 5 & 7 & 13 & 23 & 3 \\
\end{bktable}
\end{bkcenter}
\end{yattemiyou}

\noindent 1を\R{言}{い}うと5、2を\R{言}{い}うと7です。\R{言}{い}った\R{数}{かず}が1ふえると、\R{答}{こた}えは2ふえています。だから「2\R{倍}{ばい}」が\R{入}{はい}っていそうです。$1\times2=2$ で、5になるには3たりません。ほかの\R{数}{かず}でも\R{確}{たし}かめると、$5\times2+3=13$、$10\times2+3=23$、$0\times2+3=3$ で、\R{全部}{ぜん|ぶ}\R{合}{あ}います。きまりは「2\R{倍}{ばい}して3たす」、\R{式}{しき}で\R{書}{か}くと $y=2x+3$ です。

このゲームで\R{大事}{だい|じ}なのは、\textbf{\R{同}{おな}じ\R{数}{かず}を\R{言}{い}えば、いつも\R{同}{おな}じ\R{答}{こた}えが\R{返}{かえ}ってくる}ことです。もし、1を\R{言}{い}ったのに、あるときは5、あるときは8と\R{答}{こた}えが\R{変}{か}わったら、きまりは\R{当}{あ}てられません。この\R{章}{しょう}では、この「\R{入}{い}れると、\R{決}{き}まった1つが\R{出}{で}てくる」しくみに\R{名前}{な|まえ}をつけて、\R{書}{か}き\R{方}{かた}を\R{学}{まな}びます。

\begin{omoidasou}
\bkrecall{2-11}{\R{文字}{も|じ}を\R{使}{つか}うと、どんな\R{数}{かず}でも\R{成}{な}り\R{立}{た}つきまりを1つの\R{式}{しき}で\R{書}{か}ける。}
\bkrecall{2-13}{\R{文字}{も|じ}に\R{数}{かず}を\R{入}{い}れるのが\R{代入}{だい|にゅう}。マイナスの\R{数}{かず}や\R{式}{しき}は、かっこをつけて、かっこごと\R{入}{い}れる。}
\bkrecall{2-15}{\R{比例}{ひ|れい} $y=ax$：$x$ が2\R{倍}{ばい}・3\R{倍}{ばい}になると $y$ も2\R{倍}{ばい}・3\R{倍}{ばい}。グラフは\R{原点}{げん|てん}を\R{通}{とお}る\R{直線}{ちょく|せん}。}
\bkrecall{3-1}{$x$ を1つ\R{決}{き}めると $y$ がただ1つに\R{決}{き}まる\R{関係}{かん|けい}がある。まわりが\qty{20}{cm}の\R{長方形}{ちょう|ほう|けい}の\R{面積}{めん|せき}は、1つに\R{決}{き}まらない。\R{変数}{へん|すう}のとる\R{値}{あたい}の\R{範囲}{はん|い}（\R{変域}{へん|いき}）は、\R{不等号}{ふ|とう|ごう}で\R{書}{か}く。}
\bkrecall{3-2}{\R{点}{てん} $(x,\ y)$ は、\R{原点}{げん|てん}から\R{右}{みぎ}へ $x$、\R{上}{うえ}へ $y$ \R{進}{すす}んだ\R{所}{ところ}にかく。}
\end{omoidasou}

\bkfukidashi{mina}{きまりを\R{当}{あ}てるとき、\R{表}{ひょう}を\R{横}{よこ}に\R{見}{み}たら「1ふえると2ふえる」が\R{見}{み}つかったよ。}

%% ---------------- 数学の道具 ----------------
\bksection{doguu}{\R{数学}{すう|がく}の\R{道具}{どう|ぐ}}

\bksubsection{1}{1つ\R{決}{き}まると、1つ\R{決}{き}まる：\R{関数}{かん|すう}}

\leavevmode\label{bk:3-3:kansu}3-1 では、1つの\R{量}{りょう}を\R{決}{き}めると、もう1つの\R{量}{りょう}がただ1つに\R{決}{き}まる\R{関係}{かん|けい}を\R{見}{み}ました。ここでは、この\R{関係}{かん|けい}に\R{名前}{な|まえ}をつけます。

\R{身}{み}のまわりには、\R{一方}{いっ|ぽう}が\R{変}{か}わると、それにつれてもう\R{一方}{いっ|ぽう}も\R{変}{か}わる2つの\R{量}{りょう}がたくさんあります。たとえば、\R{正方形}{せい|ほう|けい}の1\R{辺}{ぺん}の\R{長}{なが}さを\R{変}{か}えると、\R{面積}{めん|せき}も\R{変}{か}わります。このような2つの\R{量}{りょう}を「\textbf{ともなって\R{変}{か}わる\R{量}{りょう}}」といいます。

\R{正方形}{せい|ほう|けい}の1\R{辺}{ぺん}を $x$〔cm〕、\R{面積}{めん|せき}を $y$〔\unit{cm^2}〕として、\R{表}{ひょう}にしてみます。\R{面積}{めん|せき}は1\R{辺}{ぺん}×1\R{辺}{ぺん}なので、$x=3$ なら $y=3\times3=9$ です。
\begin{bkcenter}
\begin{bktable}{colspec={Q[l,m,wd=38mm] *{5}{Q[c,m,wd=11mm]}}, rowsep=1.5pt}
1\R{辺}{ぺん} $x$（cm） & 1 & 2 & 3 & 4 & 5 \\
\R{面積}{めん|せき} $y$（\unit{cm^2}） & 1 & 4 & 9 & 16 & 25 \\
\end{bktable}
\bktabcaption{1\R{辺}{ぺん}を\R{決}{き}めると、\R{面積}{めん|せき}がただ1つに\R{決}{き}まる。}
\end{bkcenter}
1\R{辺}{ぺん}を3\,cmと\R{決}{き}めれば、\R{面積}{めん|せき}は\R{必}{かなら}ず\qty{9}{cm^2}です。\R{同}{おな}じ3\,cmなのに、\R{面積}{めん|せき}が\qty{9}{cm^2}になったり\qty{10}{cm^2}になったりすることはありません。

\textbf{\R{決}{き}まらない\R{例}{れい}}\quad \R{今度}{こん|ど}は、クラスの\R{人}{ひと}の「\R{年齢}{ねん|れい}」と「\R{身長}{しん|ちょう}」を\R{考}{かんが}えます。\R{同}{おな}じ13\R{歳}{さい}でも、\R{身長}{しん|ちょう}が150\,cmの\R{人}{ひと}も、162\,cmの\R{人}{ひと}もいます。\R{年齢}{ねん|れい}を13\R{歳}{さい}と\R{決}{き}めても、\R{身長}{しん|ちょう}は1つに\R{決}{き}まりません。
\begin{bkfigure}
\begin{tikzpicture}[x=1mm, y=1mm]
  % 左：自動販売機（ボタン1つに飲み物1つ）
  \fill[bkRed!75, rounded corners=2mm] (0,0) rectangle (30,44);
  \fill[white, rounded corners=1mm] (3,22) rectangle (27,41);
  \foreach \i/\c in {0/bkBlue,1/bkGreen,2/bkOrange} {
    \fill[\c!70, rounded corners=0.8mm] (5+\i*7.6,26) rectangle ++(5.2,12);
    \fill[bkInk!80, rounded corners=0.6mm] (5.6+\i*7.6,18) rectangle ++(4,2.6);}
  \fill[bkInk!80, rounded corners=1mm] (5,4) rectangle (25,10);
  \node[font=\sffamily\scriptsize, text=white] at (15,14) {ボタン};
  \node[font=\sffamily\footnotesize, align=center, anchor=north] at (15,-2) {\R{自動販売機}{じ|どう|はん|ばい|き}\\ボタン1つに、\R{飲}{の}み\R{物}{もの}1つ};
  \node[font=\sffamily\bfseries\small, text=bkGreen!50!black] at (15,52) {\R{関数}{かん|すう}（○）};
  % 右：ガチャ（\R{同}{おな}じお金で、何が出るかわからない）
  \begin{scope}[shift={(72,0)}]
  % 台（下の箱）と、丸いケースをつなぐ首
  \fill[bkBlue!70] (10,16) rectangle (20,21);
  \fill[bkBlue!70, rounded corners=2mm] (-2,0) rectangle (32,18);
  % 丸いケース（半径13）。玉は半径2.8で、中心どうしを7.2はなすので重ならない。いちばん外の玉もケースの内側に収まる
  \draw[bkBlue!70, line width=1.2pt, fill=white] (15,32) circle (13);
  \foreach \a/\c in {30/bkRed,90/bkYellow,150/bkGreen,210/bkOrange,270/bkPurple,330/bkBlue} {
    \fill[\c!65] ($(15,32)+(\a:7.2)$) circle (2.8);
    \draw[white, line width=0.5pt] ($(15,32)+(\a:7.2)+(-2.8,0)$) -- ++(5.6,0);}
  \fill[bkRed!45] (15,32) circle (2.8);
  \draw[white, line width=0.5pt] (12.2,32) -- (17.8,32);
  % お金を入れる口
  \fill[bkInk!80, rounded corners=0.5mm] (2.5,13.4) rectangle (9.5,15.2);
  % ハンドル
  \fill[white] (6,7) circle (4.2);
  \fill[bkInk!80, rotate around={30:(6,7)}] (4.6,3.8) rectangle (7.4,10.2);
  % 取り出し口と、出てきたカプセル（何色かわからない）
  \fill[bkInk!80, rounded corners=1.2mm] (16,2) rectangle (30,13);
  \fill[white] (23,7.2) circle (3.8);
  \draw[bkGray, line width=0.6pt, dash pattern=on 1.2pt off 0.9pt] (23,7.2) circle (3.8);
  \node[font=\sffamily\bfseries\small, text=bkInk, inner sep=0pt] at (23,7.2) {？};
  \node[font=\sffamily\footnotesize, align=center, anchor=north] at (15,-2) {カプセルのおもちゃ\\\R{同}{おな}じお\R{金}{かね}でも、\R{何}{なに}が\R{出}{で}るかわからない};
  \node[font=\sffamily\bfseries\small, text=bkRed!80!black] at (15,52) {\R{関数}{かん|すう}ではない（×）};
  \end{scope}
\end{tikzpicture}
\bkfigcaption{\R{入}{い}れるもの（ボタン・お\R{金}{かね}）を\R{決}{き}めると、\R{出}{で}てくるものが1つに\R{決}{き}まるか？}
\end{bkfigure}

\textbf{もう1つの\R{見方}{み|かた}：\R{自動販売機}{じ|どう|はん|ばい|き}}\quad \R{自動販売機}{じ|どう|はん|ばい|き}は、どのボタンをおすかを\R{決}{き}めると、\R{出}{で}てくる\R{飲}{の}み\R{物}{もの}が1つに\R{決}{き}まります。ちがうボタンから\R{同}{おな}じお\R{茶}{ちゃ}が\R{出}{で}てもかまいません。でも、カプセルのおもちゃの\R{機械}{き|かい}は、\R{同}{おな}じお\R{金}{かね}を\R{入}{い}れても、\R{何}{なに}が\R{出}{で}てくるかわかりません。\R{出}{で}てくるものが1つに\R{決}{き}まらないのです。

このように、2つの\R{量}{りょう} $x$ と $y$ があって、\textbf{$x$ の\R{値}{あたい}を1つ\R{決}{き}めると、それにともなって $y$ の\R{値}{あたい}がただ1つに\R{決}{き}まる}とき、「$y$ は $x$ の\textbf{\R{関数}{かん|すう}}である」といいます。\bkidx{かんすう}{\R{関数}{かん|すう}}\R{正方形}{せい|ほう|けい}の\R{面積}{めん|せき}は1\R{辺}{ぺん}の\R{関数}{かん|すう}です。\R{身長}{しん|ちょう}は\R{年齢}{ねん|れい}の\R{関数}{かん|すう}ではありません。$x$ と $y$ は、どちらも\R{値}{あたい}が\R{変}{か}わる\R{変数}{へん|すう}です（3-1）。$x$ を「\R{入}{い}れる\R{数}{かず}」、$y$ を「\R{出}{で}てくる\R{数}{かず}」と\R{考}{かんが}えると、\R{関数}{かん|すう}は「\R{入}{い}れると、\R{決}{き}まった1つが\R{出}{で}てくるルール」です。

\begin{seikaku}
「ただ1つに\R{決}{き}まる」とは、\R{同}{おな}じ $x$ に\R{対}{たい}して $y$ が2つ\R{以上}{い|じょう}にならない、ということです。ちがう $x$ に\R{対}{たい}して\R{同}{おな}じ $y$ になるのはかまいません。たとえば、$y=x^2$ では、$x=3$ でも $x=-3$ でも $y=9$ ですが、$y$ は $x$ の\R{関数}{かん|すう}です。
\end{seikaku}

\noindent\begin{minipage}{\linewidth}
\begin{motto}[2]
\textbf{グラフで\R{関数}{かん|すう}を\R{見分}{み|わ}ける：たての\R{線}{せん}を\R{引}{ひ}いてみよう}\par
\noindent\begin{minipage}[c]{0.57\linewidth}
グラフに、たての\R{直線}{ちょく|せん}（$x$ が\R{同}{おな}じ\R{値}{あたい}の\R{点}{てん}の\R{集}{あつ}まり）を、いろいろな\R{所}{ところ}に\R{引}{ひ}いてみます。どこに\R{引}{ひ}いても、グラフと\R{交}{まじ}わる\R{点}{てん}が1つまでなら、$y$ は $x$ の\R{関数}{かん|すう}です。まっすぐなグラフはそうなっています。\R{円}{えん}の\R{形}{かたち}のグラフは、たての\R{線}{せん}と2つの\R{点}{てん}で\R{交}{まじ}わる\R{所}{ところ}があります。\R{同}{おな}じ $x$ に $y$ が2つあるので、\R{関数}{かん|すう}ではありません。
\end{minipage}\hfill
\begin{minipage}[c]{0.4\linewidth}\centering
\begin{tikzpicture}[x=0.92mm, y=0.92mm]
  \draw[-{Stealth[length=3pt]}, bkInk, line width=0.6pt] (0,0) -- (20,0);
  \draw[-{Stealth[length=3pt]}, bkInk, line width=0.6pt] (2,-2) -- (2,22);
  \draw[bkBlue, line width=1.2pt] (3,4) -- (18,18);
  \draw[bkRed, line width=0.7pt, dash pattern=on 2pt off 1.5pt] (11,-1) -- (11,21);
  \fill[bkRed] (11,11.47) circle (0.9);
  \node[font=\sffamily\footnotesize, text=bkGreen!50!black, anchor=north] at (10,-2.5) {\R{関数}{かん|すう}};
  \begin{scope}[shift={(26,0)}]
  \draw[-{Stealth[length=3pt]}, bkInk, line width=0.6pt] (0,0) -- (22,0);
  \draw[-{Stealth[length=3pt]}, bkInk, line width=0.6pt] (2,-2) -- (2,22);
  \draw[bkBlue, line width=1.2pt] (12,11) circle (8);
  \draw[bkRed, line width=0.7pt, dash pattern=on 2pt off 1.5pt] (9,-1) -- (9,21);
  \fill[bkRed] (9,18.42) circle (0.9) (9,3.58) circle (0.9);
  \node[font=\sffamily\footnotesize, text=bkRed!80!black, anchor=north] at (12,-2.5) {\R{関数}{かん|すう}ではない};
  \end{scope}
\end{tikzpicture}
\end{minipage}
\end{motto}
\end{minipage}

\bksubsection{2}{3つの\R{見方}{み|かた}：\R{表}{ひょう}・グラフ・\R{式}{しき}}

\leavevmode\label{bk:3-3:mikata}\bkhatenaref の\R{温度}{おん|ど}で\R{考}{かんが}えましょう。\R{華氏}{か|し}と\R{摂氏}{せっ|し}は、\R{目}{め}もりの0の\R{場所}{ば|しょ}も、1\R{目}{め}もりの\R{大}{おお}きさもちがいます。そこで、2つの\R{目}{め}もりをくらべるのに、だれでも\R{作}{つく}れる2つの\R{温度}{おん|ど}を\R{使}{つか}います。
\begin{itemize}[leftmargin=1.6\zw, itemsep=1pt, topsep=2pt, label={\color{\bkTT}●}]
\item \R{水}{みず}がこおる\R{温度}{おん|ど}：\R{摂氏}{せっ|し}で\qty{0}{\degreeCelsius}、\R{華氏}{か|し}で\bkthrdF{32}
\item \R{水}{みず}がふっとうする\R{温度}{おん|ど}（\R{山}{やま}の\R{上}{うえ}でない\R{所}{ところ}で）：\R{摂氏}{せっ|し}で\qty{100}{\degreeCelsius}、\R{華氏}{か|し}で\bkthrdF{212}
\end{itemize}

\begin{bkfigure}
\begin{tikzpicture}[x=1mm, y=1mm]
  % 2本の温度計の目もりを横にならべる（摂氏 0〜100 と華氏 32〜212 が\R{同}{おな}じ高さ）
  \draw[bkBlue, line width=1pt] (30,0) -- (30,60);
  \draw[bkOrange, line width=1pt] (80,0) -- (80,60);
  \node[font=\sffamily\bfseries\small, text=bkBlue!80!black] at (30,66) {\R{摂氏}{せっ|し}（℃）};
  \node[font=\sffamily\bfseries\small, text=bkOrange!80!black] at (80,66) {\R{華氏}{か|し}（°F）};
  \foreach \k in {0,...,10} {\draw[bkBlue, line width=0.6pt] (27.5,\k*6) -- (30,\k*6);}
  \foreach \c in {0,20,40,60,80,100} {\node[font=\footnotesize, anchor=east] at (26.5,\c*0.6) {\c};}
  \foreach \k in {0,...,18} {\draw[bkOrange, line width=0.5pt] (80,\k*10/3) -- (82,\k*10/3);}
  \foreach \f in {32,68,104,140,176,212} {\node[font=\footnotesize, anchor=west] at (83.5,{(\f-32)/3}) {\f};}
  \draw[bkGray, line width=0.6pt, dash pattern=on 2pt off 1.5pt] (31.5,0) -- (78.5,0) (31.5,60) -- (78.5,60);
  \node[font=\sffamily\footnotesize, fill=white, inner sep=1pt] at (55,0) {\R{水}{みず}がこおる};
  \node[font=\sffamily\footnotesize, fill=white, inner sep=1pt] at (55,60) {\R{水}{みず}がふっとうする};
  \draw[{Stealth[length=3pt]}-{Stealth[length=3pt]}, bkBlue, line width=0.6pt] (12,0.8) -- (12,59.2);
  \node[font=\sffamily\footnotesize, text=bkBlue!80!black, fill=white, inner sep=1pt, align=center] at (12,30) {100\\\R{目}{め}もり};
  \draw[{Stealth[length=3pt]}-{Stealth[length=3pt]}, bkOrange, line width=0.6pt] (100,0.8) -- (100,59.2);
  \node[font=\sffamily\footnotesize, text=bkOrange!80!black, fill=white, inner sep=1pt, align=center] at (100,30) {180\\\R{目}{め}もり};
\end{tikzpicture}
\bkfigcaption{\R{同}{おな}じ\R{温度}{おん|ど}のはばが、\R{摂氏}{せっ|し}では100\R{目}{め}もり、\R{華氏}{か|し}では180\R{目}{め}もり。}
\end{bkfigure}

\noindent こおる\R{温度}{おん|ど}からふっとうする\R{温度}{おん|ど}までは、\R{摂氏}{せっ|し}では $100-0=100$ \R{目}{め}もり、\R{華氏}{か|し}では $212-32=180$ \R{目}{め}もりです。\R{同}{おな}じはばを、\R{華氏}{か|し}は180に\R{分}{わ}けているのです。だから、\R{摂氏}{せっ|し}の1\R{目}{め}もりは、\R{華氏}{か|し}では
\[ 180\div100=1.8 \]
\R{目}{め}もり\R{分}{ぶん}です（1あたりの\R{量}{りょう}。2-3）。\R{摂氏}{せっ|し}が\qty{1}{\degreeCelsius}\R{上}{あ}がるごとに、\R{華氏}{か|し}は\bkthrdF{1.8}ずつ\R{上}{あ}がります。

\bkthrdstep{1}{\R{表}{ひょう}にする}
\R{摂氏}{せっ|し}が0℃のとき、\R{華氏}{か|し}は32°F。そこから\R{摂氏}{せっ|し}が10℃\R{上}{あ}がるごとに、\R{華氏}{か|し}は $1.8\times10=18$ ずつ\R{上}{あ}がります。
\begin{bkcenter}
\begin{bktable}{colspec={Q[l,m,wd=30mm] *{7}{Q[c,m,wd=10.5mm]}}, rowsep=1.5pt}
\R{摂氏}{せっ|し} $C$（℃） & $-20$ & $-10$ & 0 & 10 & 20 & 30 & 40 \\
\R{華氏}{か|し} $F$（°F） & $-4$ & 14 & 32 & 50 & 68 & \SetCell{bg=bkYellow!40} 86 & 104 \\
\end{bktable}
\bktabcaption{\R{摂氏}{せっ|し}が10ふえるごとに、\R{華氏}{か|し}は18ふえる。}
\end{bkcenter}
\R{表}{ひょう}を\R{見}{み}ると、\bkthrdF{86}は\qty{30}{\degreeCelsius}です。\bkhatenaref の\R{答}{こた}えは\qty{30}{\degreeCelsius}。\R{暑}{あつ}い\R{夏}{なつ}の\R{日}{ひ}です。

\bkthrdstep{2}{グラフにする}
\noindent\begin{minipage}[c]{0.52\linewidth}\centering
\begin{tikzpicture}
\begin{axis}[bkplot, width=62mm, height=56mm, xmin=-25, xmax=48, ymin=-15, ymax=125,
  xtick={10,20,30,40}, ytick={20,40,60,80,100,120}, xlabel={$C$（℃）}, ylabel={$F$（°F）}]
  \addplot[bkcurve, domain=-23:44] {1.8*x+32};
  \addplot[only marks, mark=*, mark size=1.6pt, bkInk] coordinates {(-20,-4) (-10,14) (0,32) (10,50) (20,68) (30,86) (40,104)};
  \draw[bkRed, line width=0.7pt, dash pattern=on 2pt off 1.5pt] (axis cs:0,86) -- (axis cs:30,86) -- (axis cs:30,0);
  \node[font=\sffamily\footnotesize, text=bkRed!80!black, anchor=south east, inner sep=1pt] at (axis cs:28,88) {86°F};
  \node[font=\sffamily\footnotesize, text=bkBlue!80!black, anchor=north west] at (axis cs:4,24) {32};
\end{axis}
\end{tikzpicture}
\end{minipage}\hfill
\begin{minipage}[c]{0.45\linewidth}
\R{表}{ひょう}の\R{組}{くみ}を\R{点}{てん}にして（3-2）、なめらかにつなぐと、まっすぐな\R{線}{せん}になります。\R{摂氏}{せっ|し}が\R{同}{おな}じだけ\R{上}{あ}がると、\R{華氏}{か|し}も\R{同}{おな}じだけ\R{上}{あ}がるからです。

グラフなら、\R{表}{ひょう}にない\R{温度}{おん|ど}も\R{読}{よ}めます。たて\R{軸}{じく}の86から\R{横}{よこ}へ\R{進}{すす}み、\R{線}{せん}に\R{当}{あ}たった\R{所}{ところ}から\R{下}{した}へ\R{進}{すす}むと、30℃です（\R{赤}{あか}い\R{点線}{てん|せん}）。
\end{minipage}

\par\vspace{3pt}
\R{線}{せん}は、たて\R{軸}{じく}と32の\R{所}{ところ}で\R{交}{まじ}わります。\R{原点}{げん|てん}は\R{通}{とお}りません。0℃が0°Fではないからです。だから、\R{華氏}{か|し}は\R{摂氏}{せっ|し}に\R{比例}{ひ|れい}はしていません（2-15）。

\bkthrdstep{3}{\R{式}{しき}にする}
\R{摂氏}{せっ|し} $C$〔℃〕は、0℃から $C$ だけ\R{上}{あ}がった\R{温度}{おん|ど}です。\R{華氏}{か|し}は32から、1.8の $C$ \R{倍}{ばい}だけ\R{上}{あ}がります。
\[ F=1.8C+32 \]
$C=30$ を\R{代入}{だい|にゅう}すると（2-13）、$F=1.8\times30+32=54+32=86$ で、\R{表}{ひょう}と\R{合}{あ}います。\R{式}{しき}があれば、\R{表}{ひょう}にもグラフにもない\R{温度}{おん|ど}でも、\R{正確}{せい|かく}に\R{計算}{けい|さん}できます。たとえば\R{体温}{たい|おん}の\qty{36.5}{\degreeCelsius}は、$1.8\times36.5+32=65.7+32=97.7$ で、\bkthrdF{97.7}です。

\textbf{3つは\R{同}{おな}じことを\R{言}{い}っている}\quad ソラさんの\R{表}{ひょう}、ミナさんのグラフ、デルタさんの\R{式}{しき}は、どれも「\R{摂氏}{せっ|し}が\R{決}{き}まると、\R{華氏}{か|し}がただ1つに\R{決}{き}まる」という\R{同}{おな}じ\R{関数}{かん|すう}を\R{表}{あらわ}しています。\R{言}{い}い\R{方}{かた}がちがうだけです。それぞれに、\R{得意}{とく|い}なことがあります。
\begin{bkcenter}
\begin{bktable}{colspec={Q[c,m,wd=16mm] Q[l,m,wd=50mm] Q[l,m,wd=50mm]}, rowsep=1.5pt, cells={font=\small}, row{1}={font=\sffamily\bfseries\small}}
 & \R{得意}{とく|い}なこと & にがてなこと \\
\R{表}{ひょう} & \R{書}{か}いてある\R{値}{あたい}が、すぐ\R{正確}{せい|かく}に\R{読}{よ}める & \R{書}{か}いていない\R{値}{あたい}はわからない \\
グラフ & \R{全体}{ぜん|たい}の\R{形}{かたち}と、\R{変}{か}わり\R{方}{かた}がひと\R{目}{め}でわかる & \R{細}{こま}かい\R{値}{あたい}は、\R{目}{め}もりの\R{分}{ぶん}だけおおまか \\
\R{式}{しき} & どんな\R{値}{あたい}でも、\R{計算}{けい|さん}で\R{正確}{せい|かく}に\R{求}{もと}められる & \R{形}{かたち}や\R{変}{か}わり\R{方}{かた}が、すぐには\R{見}{み}えない \\
\end{bktable}
\bktabcaption{\R{表}{ひょう}・グラフ・\R{式}{しき}の、\R{得意}{とく|い}なこととにがてなこと。}
\end{bkcenter}

\begin{bkfigure}
\begin{tikzpicture}[x=1mm, y=1mm]
  \node[bk node, minimum width=16mm, minimum height=10mm] (s) at (8,0) {\R{式}{しき}};
  \node[bk node, minimum width=16mm, minimum height=10mm] (h) at (61,0) {\R{表}{ひょう}};
  \node[bk node, minimum width=18mm, minimum height=10mm] (g) at (113,0) {グラフ};
  \draw[bk link] ([yshift=2mm]s.east) -- node[bk link label, above] {\R{代入}{だい|にゅう}する（2-13）} ([yshift=2mm]h.west);
  \draw[bk link] ([yshift=-2mm]h.west) -- node[bk link label, below] {きまりを\R{見}{み}つける（3-5）} ([yshift=-2mm]s.east);
  \draw[bk link] ([yshift=2mm]h.east) -- node[bk link label, above] {\R{点}{てん}をうつ（3-2）} ([yshift=2mm]g.west);
  \draw[bk link] ([yshift=-2mm]g.west) -- node[bk link label, below] {\R{目}{め}もりを\R{読}{よ}む（3-4）} ([yshift=-2mm]h.east);
\end{tikzpicture}
\bkfigcaption{\R{表}{ひょう}・グラフ・\R{式}{しき}は、\R{行}{い}き\R{来}{き}できる。}
\end{bkfigure}

\bksubsection{3}{\R{入}{い}れると\R{出}{で}てくる\R{箱}{はこ}：$f(x)$}

\leavevmode\label{bk:3-3:fx}\R{関数}{かん|すう}は、\R{数}{かず}を\R{入}{い}れると、きまりにしたがって1つの\R{数}{かず}が\R{出}{で}てくる「\R{箱}{はこ}」のようなものです。\R{箱}{はこ}の\R{中}{なか}のきまりが「1.8\R{倍}{ばい}して32たす」なら、30を\R{入}{い}れると86が\R{出}{で}てきます。

\begin{bkfigure}
\begin{tikzpicture}[x=1mm, y=1mm]
  \bkthrdbox{(60,0)}{\R{関数}{かん|すう} $f$}{1.8\R{倍}{ばい}して\\32たす}{36}
  \node[font=\sffamily\small] (in) at (12,0) {$x$};
  \node[font=\sffamily\small] (out) at (108,0) {$f(x)$};
  \draw[bk link] (in) -- (40,0);
  \draw[bk link] (80,0) -- (out);
  \node[font=\sffamily\footnotesize, text=bkGray, anchor=south] at (12,4) {\R{入}{い}れる\R{数}{かず}};
  \node[font=\sffamily\footnotesize, text=bkGray, anchor=south] at (108,4) {\R{出}{で}てくる\R{数}{かず}};
  \node[font=\small, text=bkBlue!80!black] (in2) at (12,-17) {30};
  \node[font=\small, text=bkBlue!80!black] (out2) at (108,-17) {$f(30)=86$};
  \draw[bk link, bkBlue!60] (in2) -- (40,-17);
  \draw[bk link, bkBlue!60] (80,-17) -- (out2);
  \node[font=\sffamily\footnotesize, text=bkBlue!80!black] at (60,-17) {$1.8\times30+32$};
\end{tikzpicture}
\bkfigcaption{\R{関数}{かん|すう} $f$ は、\R{入}{い}れた\R{数}{かず} $x$ から、1つの\R{数}{かず} $f(x)$ を\R{作}{つく}って\R{出}{だ}す\R{箱}{はこ}。}
\end{bkfigure}

この\R{箱}{はこ}に $f$ という\R{名前}{な|まえ}をつけます。$f$ は、\R{関数}{かん|すう}を\R{英語}{えい|ご}で function（ファンクション）ということから\R{来}{き}た\R{文字}{も|じ}です。$f(x)$ という\R{書}{か}き\R{方}{かた}は、およそ290\R{年前}{ねん|まえ}に、スイスの\R{数学者}{すう|がく|しゃ}オイラーが\R{使}{つか}いはじめました。\R{日本}{に|ほん}では、むかし\R{関数}{かん|すう}を「\R{函数}{かん|すう}」と\R{書}{か}きました。「\R{函}{かん}」は「はこ」という\R{意味}{い|み}の\R{字}{じ}です。そして、
\begin{itemize}[leftmargin=1.6\zw, itemsep=1pt, topsep=2pt, label={\color{\bkTT}●}]
\item \R{箱}{はこ} $f$ に $x$ を\R{入}{い}れて\R{出}{で}てきた\R{値}{あたい}を、$f(x)$ と\R{書}{か}き、「\textbf{エフ・エックス}」と\R{読}{よ}みます。
\item \R{箱}{はこ} $f$ に3を\R{入}{い}れて\R{出}{で}てきた\R{値}{あたい}は、$f(3)$ と\R{書}{か}き、「\textbf{エフ・かっこ3}」と\R{読}{よ}みます。
\end{itemize}
この\R{箱}{はこ}のきまりは、$f(x)=1.8x+32$ と\R{書}{か}けます。「$x$ を\R{入}{い}れると $1.8x+32$ が\R{出}{で}てくる」という\R{意味}{い|み}です。30を\R{入}{い}れると、
\[ f(30)=1.8\times30+32=86 \]
です。「$f(30)=86$」の1\R{行}{ぎょう}で、「\R{入}{い}れた\R{数}{かず}は30」「\R{出}{で}てきた\R{数}{かず}は86」の\R{両方}{りょう|ほう}がわかります。$y=86$ と\R{書}{か}くだけでは、\R{何}{なに}を\R{入}{い}れたのかがわかりません。これが $f(x)$ の\R{書}{か}き\R{方}{かた}のよいところです。

\textbf{$f(3)$ はかけ\R{算}{ざん}ではない}\quad $f$ は「きまりの\R{名前}{な|まえ}」で、\R{数}{かず}ではありません。だから $f(3)$ は $f\times3$ ではありません。\R{駅}{えき}の「3\R{番線}{ばん|せん}」が「\R{駅}{えき}×3」でないのと\R{同}{おな}じように、$f(3)$ のかっこは「この\R{数}{かず}を\R{入}{い}れる」という\R{印}{しるし}です。2-12 では、$3(a+2)$ のように、\R{数}{かず}や\R{文字}{も|じ}とかっこのあいだのかけ\R{算}{ざん}の\R{記号}{き|ごう}を\R{省}{はぶ}きました。でも、かっこの\R{前}{まえ}が\R{関数}{かん|すう}の\R{名前}{な|まえ}のときは、かけ\R{算}{ざん}ではありません。

\textbf{\R{表}{ひょう}で\R{見}{み}ると}\quad \R{関数}{かん|すう} $f(x)=x^2$ を、\R{表}{ひょう}にしてみます。\R{上}{うえ}の\R{行}{ぎょう}が\R{入}{い}れる\R{数}{かず} $x$、\R{下}{した}の\R{行}{ぎょう}が\R{出}{で}てくる\R{数}{かず} $f(x)$ です。$f(3)$ は「\R{表}{ひょう}の $x=3$ の\R{下}{した}の\R{数}{かず}」なので、$f(3)=9$ です。$f(-3)$ も9です。
\begin{bkcenter}
\begin{bktable}{colspec={Q[c,m,wd=14mm] *{5}{Q[c,m,wd=9mm]}}, rowsep=1.5pt}
$x$ & $-3$ & $-1$ & 0 & 2 & 3 \\
$f(x)$ & 9 & 1 & 0 & 4 & \SetCell{bg=bkYellow!40} 9 \\
\end{bktable}
\end{bkcenter}
\par\vspace{4pt}
\textbf{いろいろな\R{名前}{な|まえ}の\R{関数}{かん|すう}}\quad \R{関数}{かん|すう}がいくつも\R{出}{で}てくるときは、$g(x)$、$h(x)$ のように、ほかの\R{文字}{も|じ}も\R{使}{つか}います。\R{物理}{ぶつ|り}では、\R{量}{りょう}の\R{名前}{な|まえ}をそのまま\R{関数}{かん|すう}の\R{名前}{な|まえ}に\R{使}{つか}うことがよくあります。たとえば、ボールが\R{落}{お}ちはじめてから $t$ \R{秒}{びょう}\R{後}{ご}に\R{落}{お}ちた\R{距離}{きょ|り} $x$ は $x=5t^2$（m）でした（2-13）。$x$ は\R{時刻}{じ|こく} $t$ の\R{関数}{かん|すう}なので、
\[ x(t)=5t^2 \]
と\R{書}{か}き、「\textbf{エックス・かっこティー}」と\R{読}{よ}みます。$x(2)=5\times2^2=20$ は「2\R{秒}{びょう}\R{後}{ご}には\qty{20}{m}\R{落}{お}ちている」という\R{意味}{い|み}です。\R{時刻}{じ|こく}1.1\R{秒}{びょう}なら $x(1.1)$（エックス・かっこ1.1）と\R{書}{か}きます。\R{第}{だい}5\R{巻}{かん}では、この\R{書}{か}き\R{方}{かた}で\R{速}{はや}さを\R{計算}{けい|さん}します（5-1）。


\begin{motto}[2]
\textbf{プログラムの「\R{関数}{かん|すう}」}\par
2-19 のプログラム（Python）でも、「\R{入}{い}れると\R{出}{で}てくる\R{箱}{はこ}」を\R{作}{つく}れます。プログラムでも、この\R{箱}{はこ}を\R{関数}{かん|すう}といいます。
\par\vspace{3pt}
\noindent\begin{minipage}[c]{0.5\linewidth}
\begin{bkthrdcode}
def f(x):
    return 1.8 * x + 32
print(f(30))
print(f(100))
\end{bkthrdcode}
\end{minipage}\hfill
\begin{minipage}[c]{0.4\linewidth}
\begin{bkthrdout}
86.0
212.0
\end{bkthrdout}
\end{minipage}
\par\vspace{3pt}
\texttt{def} は、\R{英語}{えい|ご}の define（\R{決}{き}める）の\R{頭}{あたま}の3\R{文字}{も|じ}で、「\R{箱}{はこ} \texttt{f} のきまりを、こう\R{決}{き}める」という\R{意味}{い|み}です。\texttt{return} は「\R{計算}{けい|さん}した\R{数}{かず}を\R{出}{だ}す」という\R{命令}{めい|れい}です。\texttt{f(30)} と\R{書}{か}くと、\texttt{x} に30が\R{入}{はい}って $1.8\times30+32$ が\R{計算}{けい|さん}され、86.0 が\R{出}{で}てきます。\R{数学}{すう|がく}の $f(30)$ と\R{同}{おな}じ\R{書}{か}き\R{方}{かた}です。
\end{motto}

\bksubsection{4}{\R{式}{しき}をかっこごと\R{入}{い}れる}

\leavevmode\label{bk:3-3:kakko}\R{箱}{はこ}には、\R{数}{かず}だけでなく、マイナスの\R{数}{かず}も、\R{文字}{も|じ}も、\R{式}{しき}も\R{入}{い}れられます。\R{入}{い}れるときは、2-13 と\R{同}{おな}じように、\textbf{$x$ の\R{所}{ところ}に、\R{入}{い}れるものをかっこごと\R{置}{お}きます}。

\begin{bkfigure}
\begin{tikzpicture}[x=1mm, y=1mm]
  \node[font=\Large] at (30,0) {$f\bigl(\,\fcolorbox{bkBlue}{bkBlue!8}{\rule{0pt}{2.2ex}\hspace{1.4em}}\,\bigr)=\bigl(\,\fcolorbox{bkBlue}{bkBlue!8}{\rule{0pt}{2.2ex}\hspace{1.4em}}\,\bigr)^2$};
  \node[font=\sffamily\footnotesize, text=bkGray, anchor=north] at (30,-8) {\R{同}{おな}じものを、\R{両方}{りょう|ほう}の□に\R{入}{い}れる};
  \node[anchor=west, font=\normalsize, align=left] at (68,9) {$f(x)=x^2$ なら};
  \node[anchor=west, font=\normalsize, align=left] at (68,1) {$f(3)=3^2=9$};
  \node[anchor=west, font=\normalsize, align=left] at (68,-7) {$f(a+h)=(a+h)^2$};
\end{tikzpicture}
\bkfigcaption{$f(x)=x^2$ は「\R{入}{い}れたものを2\R{乗}{じょう}する」きまり。\R{入}{い}れたものはかっこでつつむ。}
\end{bkfigure}

$f(x)=x^2$ なら、
\begin{itemize}[leftmargin=1.6\zw, itemsep=1pt, topsep=2pt, label={\color{\bkTT}●}]
\item 3を\R{入}{い}れると、$f(3)=3^2=9$
\item $-3$ を\R{入}{い}れると、$f(-3)=(-3)^2=(-3)\times(-3)=9$（マイナスの\R{数}{かず}はかっこごと。2-13）
\item $a+h$ を\R{入}{い}れると、$f(a+h)=(a+h)^2$
\end{itemize}
\R{最後}{さい|ご}の $f(a+h)$ は、「$a$ より $h$ だけ\R{大}{おお}きい\R{数}{かず}を\R{入}{い}れた\R{値}{あたい}」です。$a+h$ はひとまとまりの\R{数}{かず}なので、かっこでつつんでから2\R{乗}{じょう}します。$(a+h)^2$ のかっこをはずす\R{計算}{けい|さん}は、3-8 で\R{学}{まな}びます。

\textbf{\R{数}{かず}で\R{確}{たし}かめる}\quad $a=1$、$h=0.1$ とします。$a+h=1.1$ なので、はじめから1.1を\R{入}{い}れると、
\[ f(1.1)=1.1^2=1.1\times1.1=1.21 \]
です。かっこごと\R{入}{い}れた\R{式}{しき}に $a=1$、$h=0.1$ を\R{入}{い}れても、
\[ (a+h)^2=(1+0.1)^2=1.1^2=1.21 \]
で、\R{同}{おな}じになります。もし、かっこをつけずに $a+h^2$ とすると、
\[ 1+0.1^2=1+0.01=1.01 \]
で、ちがう\R{値}{あたい}になってしまいます。2\R{乗}{じょう}が $h$ だけにかかってしまうからです。

\textbf{$x$ が\R{何回}{なん|かい}\R{出}{で}てきても、\R{全部}{ぜん|ぶ}おきかえる}\quad $f(x)=3x+1$ なら、$f(a+h)$ は $x$ の\R{所}{ところ}に $a+h$ を\R{入}{い}れて、
\[ f(a+h)=3(a+h)+1=3a+3h+1 \]
です（\R{分配法則}{ぶん|ぱい|ほう|そく}。2-12）。$f(x)=x^2+x$ のように $x$ が2\R{回}{かい}\R{出}{で}てくるときは、2つとも $a+h$ におきかえて、$f(a+h)=(a+h)^2+(a+h)$ です。1つだけおきかえて、もう1つを $x$ のまま\R{残}{のこ}してはいけません。

\textbf{\R{何}{なん}のために？}\quad \R{第}{だい}5\R{巻}{かん}では、「$a$ から、ほんの\R{少}{すこ}し $h$ だけ\R{進}{すす}んだ\R{所}{ところ}」の\R{値}{あたい} $f(a+h)$ と、もとの $f(a)$ をくらべて、\R{変}{か}わる\R{速}{はや}さを\R{調}{しら}べます（5-3）。$h$ には0.1や0.01などの\R{小}{ちい}さな\R{数}{かず}が\R{入}{はい}ります。どんな $h$ でも\R{使}{つか}えるように、\R{文字}{も|じ}のまま\R{計算}{けい|さん}するのです。

\bksubsection{5}{\R{箱}{はこ}をつなぐ：\R{合成}{ごう|せい}\R{関数}{かん|すう}}

\leavevmode\label{bk:3-3:gousei}1つの\R{箱}{はこ}から\R{出}{で}てきた\R{数}{かず}を、\R{別}{べつ}の\R{箱}{はこ}に\R{入}{い}れることもできます。まず、\R{買}{か}い\R{物}{もの}で\R{考}{かんが}えましょう。1000\R{円}{えん}のおもちゃを、「100\R{円}{えん}\R{引}{び}き」のクーポンを\R{使}{つか}って\R{買}{か}います。お\R{店}{みせ}では、ねだんに\R{消費税}{しょう|ひ|ぜい}10\%がかかり、ねだんの1.1\R{倍}{ばい}をはらいます（2-5）。

\begin{bkfigure}
\begin{tikzpicture}[x=1mm, y=1mm]
  \node[font=\small] (a) at (4,0) {1000\R{円}{えん}};
  \bkthrdbox{(36,0)}{クーポン}{100\R{円}{えん}\\ひく}{24}
  \bkthrdbox{(80,0)}{\R{消費税}{しょう|ひ|ぜい}}{1.1\R{倍}{ばい}\\する}{24}
  \node[font=\small] (b) at (115,0) {990\R{円}{えん}};
  \draw[bk link] (a) -- (24,0);
  \draw[bk link] (48,0) -- node[bk link label, above] {900\R{円}{えん}} (68,0);
  \draw[bk link] (92,0) -- (b);
\end{tikzpicture}
\bkfigcaption{はじめに100\R{円}{えん}ひいて、そのあとで1.1\R{倍}{ばい}する。}
\end{bkfigure}
\noindent 2つめの\R{箱}{はこ}では、$900\times1.1=990$ になります。はらうのは990\R{円}{えん}です。

では、\R{箱}{はこ}の\R{順番}{じゅん|ばん}を\R{入}{い}れかえると、どうなるでしょう。\R{先}{さき}に1.1\R{倍}{ばい}すると $1000\times1.1=1100$ \R{円}{えん}、そのあとで100\R{円}{えん}ひくと $1100-100=1000$ \R{円}{えん}です。\textbf{\R{同}{おな}じ2つの\R{箱}{はこ}でも、つなぐ\R{順番}{じゅん|ばん}で、\R{出}{で}てくる\R{数}{かず}が\R{変}{か}わる}ことがあります。

\textbf{\R{数学}{すう|がく}の\R{箱}{はこ}でつなぐ}\quad $x$ を\R{入}{い}れると $u=3x+1$ が\R{出}{で}てくる\R{箱}{はこ}と、$u$ を\R{入}{い}れると $y=u^2$ が\R{出}{で}てくる\R{箱}{はこ}を、この\R{順}{じゅん}につなぎます。$x=0$、1、2 を\R{入}{い}れて、\R{表}{ひょう}にしてみます。

\begin{bkcenter}
\begin{bktable}{colspec={Q[l,m,wd=46mm] *{3}{Q[c,m,wd=12mm]}}, rowsep=1.5pt}
\R{入}{い}れる\R{数}{かず} $x$ & 0 & 1 & 2 \\
① 3\R{倍}{ばい}して1たす $u=3x+1$ & 1 & 4 & 7 \\
② それを2\R{乗}{じょう}する $y=u^2$ & 1 & 16 & 49 \\
\end{bktable}
\bktabcaption{2つの\R{箱}{はこ}を\R{順}{じゅん}に\R{通}{とお}す。$x=2$ なら、$u=3\times2+1=7$、$y=7^2=49$。}
\end{bkcenter}

\begin{bkfigure}
\begin{tikzpicture}[x=1mm, y=1mm]
  \node[font=\small] (x) at (2,0) {$x$};
  \bkthrdbox{(34,0)}{\R{内側}{うち|がわ}の\R{関数}{かん|すう} $g$}{$u=3x+1$}{28}
  \bkthrdbox{(82,0)}{\R{外側}{そと|がわ}の\R{関数}{かん|すう} $f$}{$y=u^2$}{28}
  \node[font=\small] (y) at (116,0) {$y$};
  \draw[bk link] (x) -- (20,0);
  \draw[bk link] (48,0) -- node[bk link label, above] {$u$} (68,0);
  \draw[bk link] (96,0) -- (y);
  \node[font=\sffamily\footnotesize, text=bkBlue!80!black] at (2,-12) {2};
  \node[font=\sffamily\footnotesize, text=bkBlue!80!black] at (58,-12) {7};
  \node[font=\sffamily\footnotesize, text=bkBlue!80!black] at (116,-12) {49};
\end{tikzpicture}
\bkfigcaption{2つの\R{関数}{かん|すう}を\R{順}{じゅん}につないだもの。\R{青}{あお}い\R{数}{かず}は $x=2$ のとき。}
\end{bkfigure}

$u=3x+1$ を $y=u^2$ の $u$ に、かっこごと\R{代入}{だい|にゅう}すると（2-13）、
\[ y=u^2=(3x+1)^2 \]
となって、$x$ から\R{直接}{ちょく|せつ} $y$ を\R{求}{もと}める\R{式}{しき}になります。$x=2$ を\R{入}{い}れると、$(3\times2+1)^2$ は $7^2=49$ になり、\R{表}{ひょう}と\R{合}{あ}います。

\textbf{\R{関数}{かん|すう}の\R{書}{か}き\R{方}{かた}で\R{書}{か}く}\quad 1つめの\R{箱}{はこ}に $g$、2つめの\R{箱}{はこ}に $f$ という\R{名前}{な|まえ}をつけると、$g(x)=3x+1$、$f(u)=u^2$ と\R{書}{か}けます。$x$ を $g$ に\R{入}{い}れて\R{出}{で}てきた $g(x)$ を、$f$ に\R{入}{い}れるので、
\[ y=f(g(x))=(3x+1)^2 \]
と\R{書}{か}きます。「エフ・ジー・エックス」と\R{読}{よ}みます。このように、\R{関数}{かん|すう}の\R{出力}{しゅつ|りょく}を\R{別}{べつ}の\R{関数}{かん|すう}に\R{入}{い}れてつくった\R{関数}{かん|すう}を、\textbf{\R{合成}{ごう|せい}\R{関数}{かん|すう}}といいます。\bkidx{ごうせいかんすう}{\R{合成}{ごう|せい}\R{関数}{かん|すう}}\R{先}{さき}にはたらく $g$ を\textbf{\R{内側}{うち|がわ}の\R{関数}{かん|すう}}、あとにはたらく $f$ を\textbf{\R{外側}{そと|がわ}の\R{関数}{かん|すう}}とよびます。$f(g(x))$ の\R{書}{か}き\R{方}{かた}でも、$g$ は\R{内側}{うち|がわ}のかっこの\R{中}{なか}にあります。

\textbf{\R{計算}{けい|さん}は、\R{内側}{うち|がわ}から}\quad $f(g(2))$ を\R{求}{もと}めるときは、かっこの\R{中}{なか}から\R{先}{さき}に\R{計算}{けい|さん}します（2-11）。
\[ g(2)=3\times2+1=7,\qquad f(g(2))=f(7)=7^2=49 \]

\textbf{\R{順番}{じゅん|ばん}を\R{入}{い}れかえると}\quad \R{先}{さき}に $f$（2\R{乗}{じょう}する）、あとで $g$（3\R{倍}{ばい}して1たす）にすると、$g(f(x))=3x^2+1$ です。$x=1$ で、$f(g(1))=f(4)=16$、$g(f(1))=g(1)=4$ で、ちがいます。\R{買}{か}い\R{物}{もの}のクーポンと\R{同}{おな}じで、\R{順番}{じゅん|ばん}が\R{大事}{だい|じ}です。

\bksubsection{6}{いろいろな\R{関数}{かん|すう}}

\label{bk:3-3:iroiro}\bkthrdstep{1}{\R{比例定数}{ひ|れい|てい|すう}がマイナスの\R{比例}{ひ|れい}}
\R{空気}{くう|き}の\R{温度}{おん|ど}は、\R{高}{たか}い\R{所}{ところ}ほど\R{低}{ひく}くなります。\R{山}{やま}に\R{登}{のぼ}ると、\qty{100}{m}\R{登}{のぼ}るごとに\R{約}{やく}\qty{0.6}{\degreeCelsius}\R{下}{さ}がります。\qty{1}{km}（\qty{1000}{m}）なら、$0.6\times10=6$ で、\R{約}{やく}\qty{6}{\degreeCelsius}\R{下}{さ}がります。\R{地面}{じ|めん}からの\R{高}{たか}さを $x$〔km〕、\R{地面}{じ|めん}とくらべた\R{温度}{おん|ど}の\R{変化}{へん|か}を $y$〔℃〕とすると、\R{下}{さ}がるので $y$ はマイナスの\R{数}{かず}です。

\noindent\begin{minipage}[c]{0.5\linewidth}
\begin{bktable}{colspec={Q[l,m,wd=20mm] *{4}{Q[c,m,wd=7.5mm]}}, rowsep=1.5pt, colsep=2pt, cells={font=\small}}
\R{高}{たか}さ $x$（km） & 0 & 1 & 2 & 3 \\
\R{変化}{へん|か} $y$（℃） & 0 & $-6$ & $-12$ & $-18$ \\
\end{bktable}
\par\vspace{4pt}
$x$ が1から2へ2\R{倍}{ばい}になると、$y$ も $-6$ から $-12$ へ2\R{倍}{ばい}です。$y\div x$ は、$(-6)\div1=-6$、$(-12)\div2=-6$ で、いつも $-6$ です。だから $y=-6x$ で、\R{比例定数}{ひ|れい|てい|すう}が $-6$ の\R{比例}{ひ|れい}です。
\end{minipage}\hfill
\begin{minipage}[c]{0.46\linewidth}\centering
\begin{tikzpicture}
\begin{axis}[bkplot, width=52mm, height=50mm, xmin=-0.4, xmax=4.2, ymin=-22, ymax=4, xtick={1,2,3}, ytick={-18,-12,-6}, xlabel={$x$（km）}, ylabel={$y$（℃）}]
  \addplot[bkcurve, domain=0:3] {-6*x};
  \addplot[only marks, mark=*, mark size=1.8pt, bkBlue] coordinates {(0,0) (3,-18)};
  \addplot[only marks, mark=*, mark size=1.5pt, bkInk] coordinates {(1,-6) (2,-12)};
  \node[font=\sffamily\footnotesize, text=bkBlue!80!black, anchor=south west] at (axis cs:2.3,-11) {$y=-6x$};
\end{axis}
\end{tikzpicture}
\end{minipage}

\par\vspace{3pt}
\R{比例定数}{ひ|れい|てい|すう}がマイナスのときも、$x$ が2\R{倍}{ばい}・3\R{倍}{ばい}になると $y$ も2\R{倍}{ばい}・3\R{倍}{ばい}になります。グラフは\R{原点}{げん|てん}を\R{通}{とお}る\R{直線}{ちょく|せん}ですが、\textbf{\R{右}{みぎ}\R{下}{さ}がり}になります。$x$ が1ふえるごとに、$y$ が6ずつへるからです。\R{高}{たか}さ\qty{3776}{m}の\R{富士山}{ふ|じ|さん}の\R{頂上}{ちょう|じょう}（\R{約}{やく}\qty{3.8}{km}）では、$-6\times3.8=-22.8$ で、\R{地面}{じ|めん}より\R{約}{やく}\qty{23}{\degreeCelsius}\R{低}{ひく}くなります。ふもとが\qty{25}{\degreeCelsius}なら、\R{頂上}{ちょう|じょう}は\R{約}{やく}\qty{2}{\degreeCelsius}です。

\bkthrdstep{2}{\R{変域}{へん|いき}：$x$ と $y$ のとる\R{値}{あたい}の\R{範囲}{はん|い}}
\R{高}{たか}さ\qty{3}{km}の\R{山}{やま}なら、$x$ は0から3までしかとりません。このような\R{変数}{へん|すう}のとる\R{値}{あたい}の\R{範囲}{はん|い}を\textbf{\R{変域}{へん|いき}}といい、\R{不等号}{ふ|とう|ごう}で $0\leqq x\leqq3$ と\R{書}{か}きます（3-1）。\bkidx{へんいき}{\R{変域}{へん|いき}}このとき $y$ の\R{変域}{へん|いき}も\R{決}{き}まります。$x=0$ で $y=0$、$x=3$ で $y=-18$ なので、$-18\leqq y\leqq0$ です。\R{小}{ちい}さいほうを\R{左}{ひだり}に\R{書}{か}きます。$x$ が\R{大}{おお}きいはしのとき、$y$ は\R{小}{ちい}さいはしになっています。\R{右}{みぎ}\R{下}{さ}がりのグラフだからです。

グラフでは、\R{範囲}{はん|い}のはしの\R{点}{てん}を●でかきます（$y=-6x$ のグラフの\R{青}{あお}い\R{点}{てん}）。●は「この\R{点}{てん}をふくむ」という\R{印}{しるし}です。はしをふくまないとき（$<$ のとき）は、○でかきます。

\bkthrdstep{3}{\R{階段}{かい|だん}の\R{関数}{かん|すう}：\R{式}{しき}で\R{書}{か}けなくても\R{関数}{かん|すう}}\label{bk:3-3:kaidan}
\R{駐車場}{ちゅう|しゃ|じょう}の\R{料金}{りょう|きん}が「30\R{分}{ぷん}までごとに200\R{円}{えん}」だとします。10\R{分}{ぷん}でも30\R{分}{ぷん}でも200\R{円}{えん}、31\R{分}{ぷん}なら400\R{円}{えん}です。とめた\R{時間}{じ|かん}を $x$ \R{分}{ふん}、\R{料金}{りょう|きん}を $y$ \R{円}{えん}とします（$0<x\leqq120$）。
\begin{bkfigure}
\begin{tikzpicture}
\begin{axis}[bkplot, width=96mm, height=54mm, xmin=0, xmax=135, ymin=0, ymax=940, xtick={30,60,90,120}, ytick={200,400,600,800}, xlabel={$x$（\R{分}{ふん}）}, ylabel={$y$（\R{円}{えん}）}]
  \foreach \k in {0,1,2,3} {\edef\tmp{\noexpand\addplot[bkcurve, line width=1.6pt] coordinates {(\k*30,{(\k+1)*200}) ({(\k+1)*30},{(\k+1)*200})};}\tmp}
  \addplot[only marks, mark=*, mark size=2pt, bkBlue] coordinates {(30,200) (60,400) (90,600) (120,800)};
  \addplot[only marks, mark=*, mark size=2pt, white, mark options={draw=bkBlue, line width=1pt}] coordinates {(0,200) (30,400) (60,600) (90,800)};
  \node[font=\sffamily\footnotesize, anchor=west, text=bkInk] at (axis cs:3,700) {●：ふくむ\quad ○：ふくまない};
\end{axis}
\end{tikzpicture}
\bkfigcaption{\R{駐車場}{ちゅう|しゃ|じょう}の\R{料金}{りょう|きん}。$30<x\leqq60$ のときは400\R{円}{えん}。}
\end{bkfigure}
このグラフは、\R{横}{よこ}にまっすぐな\R{線}{せん}が、\R{段}{だん}になって\R{上}{あ}がっていきます。ちょうど30\R{分}{ぷん}のときは200\R{円}{えん}なので、$x=30$ の\R{所}{ところ}は、\R{下}{した}の\R{段}{だん}のはしが●、\R{上}{うえ}の\R{段}{だん}のはしが○です。$y$ を1つの\R{式}{しき}で\R{書}{か}くことはできませんが、$x$ を\R{決}{き}めると $y$ はただ1つに\R{決}{き}まるので、\textbf{$y$ は $x$ の\R{関数}{かん|すう}}です。\R{関数}{かん|すう}は、\R{式}{しき}で\R{書}{か}けるものだけではありません。2-17 で\R{見}{み}た、なめらかな\R{電圧}{でん|あつ}を8ビットの\R{数}{かず}（0〜255）にする\R{階段}{かい|だん}も、このような\R{関数}{かん|すう}です。

\bkthrdstep{4}{\R{比例}{ひ|れい}\R{式}{しき}}
\R{比例}{ひ|れい}するとき、$x$ と $y$ の\R{比}{ひ}はいつも\R{同}{おな}じです。ばねに\R{同}{おな}じおもりを3\R{個}{こ}\R{下}{さ}げると\qty{6}{cm}のびるとき、5\R{個}{こ}\R{下}{さ}げたときののびを $y$〔cm〕とすると、のびとおもりの\R{数}{かず}の\R{比}{ひ}は\R{等}{ひと}しいので、
\[ 6:3=y:5 \]
と\R{書}{か}けます。このように、\R{比}{ひ}が\R{等}{ひと}しいことを\R{表}{あらわ}す\R{式}{しき}を\textbf{\R{比例}{ひ|れい}\R{式}{しき}}といいます。\R{比}{ひ} $a:b$ の\R{比}{ひ}の\R{値}{あたい}は $\dfrac{a}{b}$（$a\div b$）でした（2-5）。\R{比}{ひ}が\R{等}{ひと}しいとき、\R{比}{ひ}の\R{値}{あたい}も\R{等}{ひと}しいので、
\[ \frac{6}{3}=\frac{y}{5}\qquad\text{すなわち}\qquad 2=\frac{y}{5} \]
$y$ を5でわると2なので、$y=2\times5=10$ です（\R{両辺}{りょう|へん}に5をかける。2-14）。のびは\qty{10}{cm}です。


\begin{taisetsu}
$x$ を1つ\R{決}{き}めると $y$ がただ1つに\R{決}{き}まるとき、$y$ は $x$ の\R{関数}{かん|すう}。\R{表}{ひょう}・グラフ・\R{式}{しき}は、\R{同}{おな}じ\R{関数}{かん|すう}の3つの\R{見方}{み|かた}。\par\vspace{2pt}
$f(x)$ は「\R{関数}{かん|すう} $f$ に $x$ を\R{入}{い}れた\R{値}{あたい}」。$f(x)=x^2$ なら $f(3)=3^2=9$、$f(a+h)=(a+h)^2$。\R{入}{い}れるものは、かっこごと。\par\vspace{2pt}
\R{関数}{かん|すう}の\R{出力}{しゅつ|りょく}を、\R{別}{べつ}の\R{関数}{かん|すう}に\R{入}{い}れることができる（\R{合成}{ごう|せい}\R{関数}{かん|すう}）。$u=3x+1$、$y=u^2$ なら $y=(3x+1)^2$。\R{計算}{けい|さん}は\R{内側}{うち|がわ}から。
\end{taisetsu}

\begin{kotoba}
\small
\bkword{\R{関数}{かん|すう}}{$x$ の\R{値}{あたい}を1つ\R{決}{き}めると、$y$ の\R{値}{あたい}がただ1つに\R{決}{き}まるとき、「$y$ は $x$ の\R{関数}{かん|すう}」という。}
\bkword{$f(x)$}{\R{関数}{かん|すう} $f$ に $x$ を\R{入}{い}れた\R{値}{あたい}。「エフ・エックス」。$f(3)$ は「エフ・かっこ3」で、かけ\R{算}{ざん}ではない。}
\bkword{\R{合成}{ごう|せい}\R{関数}{かん|すう}}{\R{関数}{かん|すう}の\R{出力}{しゅつ|りょく}を、\R{別}{べつ}の\R{関数}{かん|すう}に\R{入}{い}れてつくった\R{関数}{かん|すう}。$f(g(x))$。\R{先}{さき}にはたらく $g$ が\R{内側}{うち|がわ}の\R{関数}{かん|すう}、あとの $f$ が\R{外側}{そと|がわ}の\R{関数}{かん|すう}。}
\bkword{\R{変域}{へん|いき}}{\R{変数}{へん|すう}のとる\R{値}{あたい}の\R{範囲}{はん|い}。$0\leqq x\leqq3$ のように\R{不等号}{ふ|とう|ごう}で\R{書}{か}く。}
\bkword{\R{比例}{ひ|れい}\R{式}{しき}}{$6:3=y:5$ のように、\R{比}{ひ}が\R{等}{ひと}しいことを\R{表}{あらわ}す\R{式}{しき}。\R{比}{ひ}の\R{値}{あたい}が\R{等}{ひと}しいことを\R{使}{つか}って\R{解}{と}く。}
\end{kotoba}

\bkkigocards{\begin{minipage}[c]{67mm}\bkkigocard{$f(x)$}{エフ・エックス}{\R{関数}{かん|すう} $f$ に $x$ を\R{入}{い}れた\R{値}{あたい}}{function}{3-3}\end{minipage}}
{\begin{minipage}[c]{52mm}\small $f(3)$ は「エフ・かっこ3」。かっこは「この\R{数}{かず}を\R{入}{い}れる」という\R{印}{しるし}で、かけ\R{算}{ざん}ではありません。$x(t)$ は「エックス・かっこティー」。\end{minipage}}




%% ---------------- 例題 ----------------
\bksection{reidai}{\R{例題}{れい|だい}でたしかめよう}

\begin{reidai}[\R{関数}{かん|すう}かどうかを\R{見分}{み|わ}ける]
$y$ は $x$ の\R{関数}{かん|すう}ですか。\R{関数}{かん|すう}でないものは、$y$ が1つに\R{決}{き}まらない\R{例}{れい}を1つあげましょう。\par
（1）1\R{本}{ぽん}60\R{円}{えん}のえんぴつを $x$ \R{本}{ほん}\R{買}{か}ったときの\R{代金}{だい|きん} $y$ \R{円}{えん}\par
（2）1\R{辺}{ぺん}が $x$〔cm〕の\R{正方形}{せい|ほう|けい}の\R{面積}{めん|せき} $y$〔\unit{cm^2}〕\par
（3）まわりの\R{長}{なが}さが $x$〔cm〕の\R{長方形}{ちょう|ほう|けい}の\R{面積}{めん|せき} $y$〔\unit{cm^2}〕\par
\kangaekata $x$ を1つ\R{決}{き}めて、$y$ が1つに\R{決}{き}まるかを\R{調}{しら}べます。\R{決}{き}まらないときは、\R{同}{おな}じ $x$ で $y$ がちがう\R{例}{れい}をあげます。\par
\bkthrdSi（1）\nobreak \R{代金}{だい|きん}は $y=60x$。$x=3$ なら $y=180$ だけ。
\bkthrdI（2）\nobreak \R{面積}{めん|せき}は1\R{辺}{ぺん}×1\R{辺}{ぺん}で $y=x^2$。$y$ は1つに\R{決}{き}まる。
\bkthrdI（3）\nobreak まわりが20\,cmでも、たて5\,cm・\R{横}{よこ}5\,cmなら\R{面積}{めん|せき}は25、
\bkthrdJ たて1\,cm・\R{横}{よこ}9\,cmなら\R{面積}{めん|せき}は9。\R{同}{おな}じ $x=20$ で、$y$ が2つ。\par
\bkthrdK（1）\nobreak \R{関数}{かん|すう}\quad（2）\nobreak \R{関数}{かん|すう}\quad（3）\nobreak \R{関数}{かん|すう}ではない（\R{例}{れい}：$x=20$ のとき、$y=25$ にも $y=9$ にもなる）
\end{reidai}

\begin{reidai}[ことばから\R{式}{しき}・\R{表}{ひょう}・グラフを\R{作}{つく}る]
ある\R{電気料金}{でん|き|りょう|きん}は、\R{使}{つか}った\R{量}{りょう}に\R{関係}{かん|けい}なくかかる\R{基本料金}{き|ほん|りょう|きん}が1000\R{円}{えん}で、さらに\R{電気}{でん|き}\qty{1}{kWh}（キロワット\R{時}{じ}）ごとに30\R{円}{えん}かかります。1か\R{月}{げつ}に $x$〔kWh〕\R{使}{つか}ったときの\R{料金}{りょう|きん}を $y$ \R{円}{えん}とします。（1）$y$ を $x$ の\R{式}{しき}で\R{表}{あらわ}しましょう。（2）$x=0$、100、200、300 のときの\R{表}{ひょう}とグラフを\R{作}{つく}りましょう。（3）\R{料金}{りょう|きん}が7000\R{円}{えん}のとき、\R{使}{つか}った\R{電気}{でん|き}は\R{何}{なん}kWhですか。\par
\kangaekata \R{料金}{りょう|きん}＝\R{基本料金}{き|ほん|りょう|きん}＋（1\,kWhあたりの\R{料金}{りょう|きん}）×（\R{使}{つか}った\R{量}{りょう}）です。（3）は、\R{式}{しき}の $y$ に7000を\R{入}{い}れて、$x$ についての\R{方程式}{ほう|てい|しき}を\R{解}{と}きます（2-16）。\par
\bkthrdSi（1）\nobreak $y=1000+30\times x=30x+1000$
\bkthrdI（2）\nobreak $x=100$ なら $y=30\times100+1000=4000$
\bkthrdJ ほかの $x$ も\R{同}{おな}じように\R{代入}{だい|にゅう}して、\R{表}{ひょう}にする。\R{表}{ひょう}の\R{組}{くみ}を\R{点}{てん}にしてグラフにする。
\par\vspace{2pt}
\noindent\begin{minipage}[c]{0.55\linewidth}\centering
\begin{bktable}{colspec={Q[l,m,wd=12mm] *{4}{Q[c,m,wd=10mm]}}, rowsep=1pt, colsep=2pt, cells={font=\small}}
$x$（kWh） & 0 & 100 & 200 & 300 \\
$y$（\R{円}{えん}） & 1000 & 4000 & 7000 & 10000 \\
\end{bktable}
\end{minipage}\hfill
\begin{minipage}[c]{0.42\linewidth}\centering
\begin{tikzpicture}
\begin{axis}[bkplot, width=48mm, height=40mm, xmin=0, xmax=345, ymin=0, ymax=11800, xtick={100,200,300}, ytick={2000,4000,6000,8000,10000}, xlabel={$x$}, ylabel={$y$}]
  \addplot[bkcurve, domain=0:320] {30*x+1000};
  \addplot[only marks, mark=*, mark size=1.5pt, bkInk] coordinates {(0,1000) (100,4000) (200,7000) (300,10000)};
\end{axis}
\end{tikzpicture}
\end{minipage}
\bkthrdI（3）\nobreak $30x+1000=7000$
\bkthrdJ $30x=7000-1000=6000$\quad 1000を\R{移項}{い|こう}する（2-14）
\bkthrdJ $x=6000\div30=200$\quad \R{両辺}{りょう|へん}を30でわる\par
\bkthrdK（1）\nobreak $y=30x+1000$\quad（2）\nobreak しきの\R{中}{なか}の\R{表}{ひょう}とグラフ\quad（3）\nobreak \qty{200}{kWh}（\R{確}{たし}かめ：$30\times200+1000=7000$）
\end{reidai}

\begingroup\tcbset{equal height group=bkthrdjA}%
\noindent\begin{minipage}[t]{0.49\linewidth}
\begin{jibun}
$y$ は $x$ の\R{関数}{かん|すう}ですか。\R{関数}{かん|すう}でないものは、$y$ が1つに\R{決}{き}まらない\R{例}{れい}を1つあげましょう。（1）\R{時速}{じ|そく}\qty{40}{km}で $x$ \R{時間}{じ|かん}\R{走}{はし}ったときの\R{道}{みち}のり $y$〔km〕（2）1\R{辺}{ぺん}が $x$〔cm〕の\R{立方体}{りっ|ぽう|たい}の\R{体積}{たい|せき} $y$〔\unit{cm^3}〕（3）\R{面積}{めん|せき}が $x$〔\unit{cm^2}〕の\R{長方形}{ちょう|ほう|けい}の、たての\R{長}{なが}さ $y$〔cm〕
\bkjans{（1）\nobreak \R{関数}{かん|すう}（$y=40x$）\quad（2）\nobreak \R{関数}{かん|すう}（$y=x^3$）\quad（3）\nobreak \R{関数}{かん|すう}ではない（\R{例}{れい}：$x=12$ のとき、たて3\,cm・\R{横}{よこ}4\,cmなら $y=3$、たて2\,cm・\R{横}{よこ}6\,cmなら $y=2$）}
\end{jibun}
\end{minipage}\hfill
\begin{minipage}[t]{0.49\linewidth}
\begin{jibun}
あるスポーツジムは、\R{入会金}{にゅう|かい|きん}が2000\R{円}{えん}で、さらに1\R{回}{かい}\R{使}{つか}うごとに500\R{円}{えん}かかります。$x$ \R{回}{かい}\R{使}{つか}ったときの\R{合計}{ごう|けい}を $y$ \R{円}{えん}とします。（1）$y$ を $x$ の\R{式}{しき}で\R{表}{あらわ}しましょう。（2）$x=0$、1、2、3、4 のときの\R{表}{ひょう}とグラフを\R{作}{つく}りましょう（ノートに）。（3）\R{合計}{ごう|けい}が6000\R{円}{えん}のとき、\R{何回}{なん|かい}\R{使}{つか}いましたか。
\bkjans{（1）\nobreak $y=500x+2000$\quad（2）\nobreak $y$ は\R{順}{じゅん}に2000、2500、3000、3500、4000。グラフは、この5つの\R{点}{てん}がならぶまっすぐな\R{線}{せん}\quad（3）\nobreak $500x+2000=6000$、$500x=4000$、$x=8$ で8\R{回}{かい}}
\end{jibun}
\end{minipage}\par
\endgroup

\begin{reidai}[$f(x)$ の\R{値}{あたい}を\R{求}{もと}める]
（1）$f(x)=1.8x+32$ のとき、$f(20)$、$f(-10)$、$f(37)$ を\R{求}{もと}めましょう。\par
（2）$x(t)=5t^2$ のとき、$x(3)$ を\R{求}{もと}めましょう。\par
\kangaekata $f(20)$ は「$x$ に20を\R{入}{い}れた\R{値}{あたい}」です。マイナスの\R{数}{かず}は、かっこをつけて\R{入}{い}れます。（2）は、\R{関数}{かん|すう}の\R{名前}{な|まえ}が $x$、\R{入}{い}れる\R{文字}{も|じ}が $t$ です。$t$ に3を\R{入}{い}れます。\par
\bkthrdSi（1）\nobreak $f(20)=1.8\times20+32=36+32=68$
\bkthrdJ $f(-10)=1.8\times(-10)+32=-18+32=14$
\bkthrdJ $f(37)=1.8\times37+32=66.6+32=98.6$
\bkthrdI（2）\nobreak $x(3)=5\times3^2=5\times9=45$\par
\bkthrdK（1）\nobreak $f(20)=68$、$f(-10)=14$、$f(37)=98.6$\quad（2）\nobreak $x(3)=45$（3\R{秒}{びょう}で\qty{45}{m}\R{落}{お}ちる）
\end{reidai}

\begin{reidai}[\R{式}{しき}をかっこごと\R{入}{い}れる]
$f(x)=x^2$、$g(x)=3x+1$ のとき、\R{次}{つぎ}の\R{式}{しき}を\R{求}{もと}めましょう。\par
（1）$f(-4)$\quad（2）$f(a+h)$\quad（3）$f(2a)$\quad（4）$g(a+h)$\par
\kangaekata $x$ の\R{所}{ところ}に、\R{入}{い}れるものをかっこごと\R{置}{お}きます。（3）の $(2a)^2$ は $2a\times2a$ です。（4）はかっこの\R{前}{まえ}に3があるので、\R{分配法則}{ぶん|ぱい|ほう|そく}でかっこをはずします（2-12）。\par
\bkthrdSi（1）\nobreak $f(-4)=(-4)^2=(-4)\times(-4)=16$
\bkthrdI（2）\nobreak $f(a+h)=(a+h)^2$
\bkthrdI（3）\nobreak $f(2a)=(2a)^2=2a\times2a=2\times2\times a\times a=4a^2$
\bkthrdI（4）\nobreak $g(a+h)=3(a+h)+1=3a+3h+1$\par
\bkthrdK（1）\nobreak 16\quad（2）\nobreak $(a+h)^2$\quad（3）\nobreak $4a^2$\quad（4）\nobreak $3a+3h+1$
\end{reidai}

\begingroup\tcbset{equal height group=bkthrdjB}%
\noindent\begin{minipage}[t]{0.49\linewidth}
\begin{jibun}
（1）$g(x)=2.5x+10$ のとき、$g(4)$、$g(-2)$、$g(0.4)$ を\R{求}{もと}めましょう。（2）$x(t)=5t^2$ のとき、$x(4)$ を\R{求}{もと}めましょう。
\bkjans{（1）\nobreak $g(4)=2.5\times4+10=20$、$g(-2)=2.5\times(-2)+10=5$、$g(0.4)=2.5\times0.4+10=11$\quad（2）\nobreak $x(4)=5\times16=80$}
\end{jibun}
\end{minipage}\hfill
\begin{minipage}[t]{0.49\linewidth}
\begin{jibun}
$f(x)=x^2$、$g(x)=2x-5$ のとき、\R{次}{つぎ}の\R{式}{しき}を\R{求}{もと}めましょう。（1）$f(-5)$\quad（2）$f(t+1)$\quad（3）$f(3a)$\quad（4）$g(a+h)$
\bkjans{（1）\nobreak 25\quad（2）\nobreak $(t+1)^2$\quad（3）\nobreak $(3a)^2=9a^2$\quad（4）\nobreak $2(a+h)-5=2a+2h-5$}
\end{jibun}
\end{minipage}\par
\endgroup

\begin{reidai}[\R{箱}{はこ}をつなぐ：\R{合成}{ごう|せい}\R{関数}{かん|すう}]
$g(x)=2x-1$、$f(u)=u^2$ とします。\par
（1）$f(g(3))$ を\R{求}{もと}めましょう。\quad（2）$y=f(g(x))$ を $x$ の\R{式}{しき}で\R{表}{あらわ}しましょう。\par
（3）$g(f(3))$ を\R{求}{もと}め、（1）とくらべましょう。\par
\kangaekata \R{計算}{けい|さん}は\R{内側}{うち|がわ}から。$f(g(3))$ は、まず $g(3)$ を\R{計算}{けい|さん}し、\R{出}{で}てきた\R{数}{かず}を $f$ に\R{入}{い}れます。（2）は、$g(x)$ をそのまま $f$ にかっこごと\R{入}{い}れます。\par
\bkthrdSi（1）\nobreak $g(3)=2\times3-1=5$、$f(g(3))=f(5)=5^2=25$
\bkthrdI（2）\nobreak $f(g(x))=f(2x-1)=(2x-1)^2$
\bkthrdJ \R{確}{たし}かめ：$x=3$ を\R{入}{い}れると $(2\times3-1)^2=5^2=25$。（1）と\R{同}{おな}じ。
\bkthrdI（3）\nobreak $f(3)=3^2=9$、$g(f(3))=g(9)=2\times9-1=17$\par
\bkthrdK（1）\nobreak 25\quad（2）\nobreak $y=(2x-1)^2$（\R{内側}{うち|がわ}の\R{関数}{かん|すう}が $g$、\R{外側}{そと|がわ}の\R{関数}{かん|すう}が $f$）\quad（3）\nobreak 17。（1）の25とちがう。
\end{reidai}

\begin{reidai}[\R{比例定数}{ひ|れい|てい|すう}がマイナスの\R{比例}{ひ|れい}と\R{変域}{へん|いき}]
$y=-2x$ で、$x$ の\R{変域}{へん|いき}が $-1\leqq x\leqq3$ のとき、$y$ の\R{変域}{へん|いき}を\R{求}{もと}めましょう。\par
\kangaekata $x$ の\R{変域}{へん|いき}の\R{両}{りょう}はしで、$y$ を\R{計算}{けい|さん}します。\R{比例定数}{ひ|れい|てい|すう}がマイナスなので、グラフは\R{右}{みぎ}\R{下}{さ}がりです。$x$ が\R{大}{おお}きいほど $y$ は\R{小}{ちい}さくなります。\par
\noindent\begin{minipage}[c]{0.6\linewidth}
\bkthrdS $x=-1$ のとき、$y=-2\times(-1)=2$
\bkthrdC $x=3$ のとき、$y=-2\times3=-6$
\bkthrdC \R{小}{ちい}さいほうの $-6$ を\R{左}{ひだり}に\R{書}{か}く。\par
\bkthrdK $-6\leqq y\leqq2$
\end{minipage}\hfill
\begin{minipage}[c]{0.36\linewidth}\centering
\begin{tikzpicture}
\begin{axis}[bkplot, width=44mm, height=48mm, xmin=-2, xmax=4.2, ymin=-7.8, ymax=4, xtick={-1,3}, ytick={-6,2}, xlabel={$x$}, ylabel={$y$}]
  \addplot[bkthin, dashed, domain=-1.8:3.8] {-2*x};
  \addplot[bkcurve, domain=-1:3] {-2*x};
  \addplot[only marks, mark=*, mark size=1.8pt, bkBlue] coordinates {(-1,2) (3,-6)};
\end{axis}
\end{tikzpicture}
\end{minipage}
\end{reidai}

\begingroup\tcbset{equal height group=bkthrdjC}%
\noindent\begin{minipage}[t]{0.49\linewidth}
\begin{jibun}
$g(x)=x+3$、$f(u)=u^2$ とします。（1）$f(g(2))$ を\R{求}{もと}めましょう。（2）$y=f(g(x))$ を $x$ の\R{式}{しき}で\R{表}{あらわ}しましょう。（3）$g(f(2))$ を\R{求}{もと}め、（1）とくらべましょう。
\bkjans{（1）\nobreak $g(2)=5$、$f(5)=25$\quad（2）\nobreak $y=(x+3)^2$\quad（3）\nobreak $f(2)=4$、$g(4)=7$。（1）とちがう}
\end{jibun}
\end{minipage}\hfill
\begin{minipage}[t]{0.49\linewidth}
\begin{jibun}
$y=-3x$ で、$x$ の\R{変域}{へん|いき}が $-1\leqq x\leqq2$ のとき、$y$ の\R{変域}{へん|いき}を\R{求}{もと}めましょう。
\bkjans{$x=-1$ で $y=3$、$x=2$ で $y=-6$。$-6\leqq y\leqq3$}
\end{jibun}
\end{minipage}\par
\endgroup

\begin{reidai}[\R{階段}{かい|だん}の\R{関数}{かん|すう}を\R{読}{よ}む]
\R{駐車場}{ちゅう|しゃ|じょう}の\R{料金}{りょう|きん}は「30\R{分}{ぷん}までごとに200\R{円}{えん}」です（\pageref{bk:3-3:kaidan}ページのグラフ）。とめた\R{時間}{じ|かん}を $x$ \R{分}{ふん}、\R{料金}{りょう|きん}を $y$ \R{円}{えん}とします。（1）45\R{分}{ふん}とめたときの\R{料金}{りょう|きん}は\R{何円}{なん|えん}ですか。（2）ちょうど60\R{分}{ぷん}とめたときの\R{料金}{りょう|きん}は\R{何円}{なん|えん}ですか。（3）\R{料金}{りょう|きん}が600\R{円}{えん}になるのは、$x$ がどんな\R{範囲}{はん|い}のときですか。\R{不等号}{ふ|とう|ごう}で\R{答}{こた}えましょう。\par
\kangaekata 30\R{分}{ぷん}ごとに\R{区切}{く|ぎ}って\R{考}{かんが}えます。\R{区切}{く|ぎ}りの\R{時間}{じ|かん}ちょうどのときは、グラフの●の\R{段}{だん}（\R{下}{した}の\R{段}{だん}）です。\par
\bkthrdSi（1）\nobreak 45\R{分}{ふん}は、30\R{分}{ぷん}より\R{長}{なが}く60\R{分}{ぷん}まで、の2つめの\R{段}{だん}。$200\times2=400$
\bkthrdI（2）\nobreak 60\R{分}{ぷん}ちょうどは、「60\R{分}{ぷん}まで」に\R{入}{はい}るので、2つめの\R{段}{だん}で400\R{円}{えん}
\bkthrdI（3）\nobreak 600\R{円}{えん}は3つめの\R{段}{だん}。60\R{分}{ぷん}より\R{長}{なが}く、90\R{分}{ぷん}まで\par
\bkthrdK（1）\nobreak 400\R{円}{えん}\quad（2）\nobreak 400\R{円}{えん}\quad（3）\nobreak $60<x\leqq90$
\end{reidai}

\begin{reidai}[\R{比例}{ひ|れい}\R{式}{しき}を\R{解}{と}く]
（1）$x:12=3:4$ の $x$ を\R{求}{もと}めましょう。\par
（2）\R{同}{おな}じ\R{紙}{かみ}20まいの\R{重}{おも}さは\qty{90}{g}でした。この\R{紙}{かみ}50まいの\R{重}{おも}さを $x$〔g〕として\R{比例}{ひ|れい}\R{式}{しき}を\R{作}{つく}り、$x$ を\R{求}{もと}めましょう。\par
\kangaekata \R{比}{ひ}が\R{等}{ひと}しいとき、\R{比}{ひ}の\R{値}{あたい}（\R{前}{まえ}の\R{数}{かず}÷\R{後}{あと}の\R{数}{かず}）も\R{等}{ひと}しくなります。（2）は、\R{重}{おも}さとまい\R{数}{かず}の\R{比}{ひ}が\R{等}{ひと}しいので、$x:50=90:20$ です。\par
\bkthrdSi（1）\nobreak $\dfrac{x}{12}=\dfrac{3}{4}$\quad \R{比}{ひ}の\R{値}{あたい}が\R{等}{ひと}しい
\bkthrdJ $x=\dfrac{3}{4}\times12=9$\quad \R{両辺}{りょう|へん}に12をかける（2-14）
\bkthrdI（2）\nobreak $x:50=90:20$ より、$\dfrac{x}{50}=\dfrac{90}{20}=4.5$
\bkthrdJ $x=4.5\times50=225$\quad \R{両辺}{りょう|へん}に50をかける\par
\bkthrdK（1）\nobreak $x=9$（\R{確}{たし}かめ：$9\div12=0.75$、$3\div4=0.75$）\quad（2）\nobreak \qty{225}{g}
\end{reidai}

\begingroup\tcbset{equal height group=bkthrdjD}%
\noindent\begin{minipage}[t]{0.49\linewidth}
\begin{jibun}
\R{貸}{か}し\R{自転車}{じ|てん|しゃ}の\R{料金}{りょう|きん}は「1\R{時間}{じ|かん}までごとに300\R{円}{えん}」です。\R{借}{か}りた\R{時間}{じ|かん}を $x$ \R{時間}{じ|かん}、\R{料金}{りょう|きん}を $y$ \R{円}{えん}とします（$0<x\leqq4$）。（1）1.5\R{時間}{じ|かん}\R{借}{か}りたときの\R{料金}{りょう|きん}は\R{何円}{なん|えん}ですか。（2）ちょうど3\R{時間}{じ|かん}\R{借}{か}りたときの\R{料金}{りょう|きん}は\R{何円}{なん|えん}ですか。（3）\R{料金}{りょう|きん}が600\R{円}{えん}になるのは、$x$ がどんな\R{範囲}{はん|い}のときですか。
\bkjans{（1）\nobreak 600\R{円}{えん}\quad（2）\nobreak 900\R{円}{えん}\quad（3）\nobreak $1<x\leqq2$}
\end{jibun}
\end{minipage}\hfill
\begin{minipage}[t]{0.49\linewidth}
\begin{jibun}
（1）$x:15=2:5$ の $x$ を\R{求}{もと}めましょう。（2）\R{同}{おな}じくぎ30\R{本}{ぽん}の\R{重}{おも}さは\qty{72}{g}でした。このくぎ80\R{本}{ぽん}の\R{重}{おも}さを $x$〔g〕として\R{比例}{ひ|れい}\R{式}{しき}を\R{作}{つく}り、$x$ を\R{求}{もと}めましょう。
\bkjans{（1）\nobreak $\dfrac{x}{15}=\dfrac{2}{5}$、$x=6$\quad（2）\nobreak $x:80=72:30$、$\dfrac{x}{80}=2.4$、$x=192$ で\qty{192}{g}}
\end{jibun}
\end{minipage}\par
\endgroup

%% ---------------- まちがえやすいところ・物理工学 ----------------
\bksection{tsumazuki}{まちがえやすいところをチェック}\label{bk:3-3:tsuma}

\begin{tsumazuki}{$f(3)$ をかけ\R{算}{ざん}と\R{読}{よ}む}
\bkng{$f(x)=x^2$ のとき、$f(3)=f\times3$}
\bkok{$f(3)$ は「$f$ に3を\R{入}{い}れた\R{値}{あたい}」。$f(3)=3^2=9$}
$f$ はきまりの\R{名前}{な|まえ}で、\R{数}{かず}ではありません。かっこは「この\R{数}{かず}を\R{入}{い}れる」という\R{印}{しるし}です。
\end{tsumazuki}

\begin{tsumazuki}{ちがう $x$ で\R{同}{おな}じ $y$ なら\R{関数}{かん|すう}ではない？}
\bkng{$y=x^2$ は、$x=3$ でも $x=-3$ でも $y=9$ なので、\R{関数}{かん|すう}ではない}
\bkok{$y=x^2$ は\R{関数}{かん|すう}。\R{同}{おな}じ $x$ に $y$ が2つあるときだけ、\R{関数}{かん|すう}ではない}
ちがうボタンから\R{同}{おな}じお\R{茶}{ちゃ}が\R{出}{で}る\R{自動販売機}{じ|どう|はん|ばい|き}と\R{同}{おな}じです。$x$ を\R{決}{き}めたとき $y$ が1つなら\R{関数}{かん|すう}です。
\end{tsumazuki}

\begin{tsumazuki}{かっこをつけずに\R{入}{い}れる}
\bkng{$f(x)=x^2$ のとき、$f(a+h)=a+h^2$}
\bkok{$f(a+h)=(a+h)^2$}
$a=1$、$h=0.1$ で\R{確}{たし}かめると、$f(1.1)=1.21$。$a+h^2=1.01$ は\R{合}{あ}いません。$a+h$ はひとまとまりで2\R{乗}{じょう}します。
\end{tsumazuki}


\begin{tsumazuki}{\R{変域}{へん|いき}で、\R{大小}{だい|しょう}を\R{逆}{ぎゃく}に\R{書}{か}く}
$y=-2x$、$1\leqq x\leqq3$ のとき
\bkng{$-2\leqq y\leqq-6$}
\bkok{$-6\leqq y\leqq-2$}
$-6$ は $-2$ より\R{小}{ちい}さい\R{数}{かず}です（2-2）。\R{小}{ちい}さいほうを\R{左}{ひだり}に\R{書}{か}きます。
\end{tsumazuki}

\begin{tsumazuki}{\R{階段}{かい|だん}のはしの\R{値}{あたい}}
\bkng{30\R{分}{ぷん}までごとに200\R{円}{えん}の\R{駐車場}{ちゅう|しゃ|じょう}に、ちょうど30\R{分}{ぷん}とめると400\R{円}{えん}}
\bkok{ちょうど30\R{分}{ぷん}は「30\R{分}{ぷん}まで」に\R{入}{はい}るので200\R{円}{えん}}
グラフのはしの●と○を\R{見}{み}て、どちらの\R{段}{だん}に\R{入}{はい}るかを\R{確}{たし}かめよう。
\end{tsumazuki}

\begin{tsumazuki}{\R{合成}{ごう|せい}\R{関数}{かん|すう}の\R{順番}{じゅん|ばん}を\R{逆}{ぎゃく}にする}
$g(x)=3x+1$、$f(u)=u^2$ のとき、$f(g(1))$ は？
\bkng{\R{先}{さき}に2\R{乗}{じょう}して $1^2=1$、それから $3\times1+1=4$}
\bkok{\R{内側}{うち|がわ}の $g$ が\R{先}{さき}。$g(1)=4$、$f(4)=16$}
かっこの\R{中}{なか}（\R{内側}{うち|がわ}）から\R{計算}{けい|さん}します。\R{順番}{じゅん|ばん}を\R{入}{い}れかえると、ちがう\R{関数}{かん|すう}になります。
\end{tsumazuki}





\begin{butsuri}
\par\addvspace{3pt}\textbf{\R{温度}{おん|ど}センサー：\R{温度}{おん|ど}を\R{入}{い}れると\R{電圧}{でん|あつ}が\R{出}{で}てくる\R{箱}{はこ}}\par\noindent
エアコンや\R{冷蔵庫}{れい|ぞう|こ}には、\R{温度}{おん|ど}をはかるセンサーが\R{入}{はい}っています。ある\R{温度}{おん|ど}センサーは、\qty{1}{\degreeCelsius}につき\qty{0.01}{V}の\R{電圧}{でん|あつ}を\R{出}{だ}します。\R{温度}{おん|ど} $T$〔℃〕を\R{入}{い}れると、\R{電圧}{でん|あつ} $V=0.01T$（V）が\R{出}{で}てくる\R{関数}{かん|すう}です。\qty{25}{\degreeCelsius}なら $0.01\times25=0.25$ で、\qty{0.25}{V}です。

\par\addvspace{3pt}\textbf{\R{箱}{はこ}をつなぐと、\R{温度}{おん|ど}が\R{数}{かず}になる}\par\noindent
コンピュータは、\R{電圧}{でん|あつ}をそのままでは\R{計算}{けい|さん}できません。\R{電圧}{でん|あつ}を\R{数}{かず}に\R{変}{か}える\R{部品}{ぶ|ひん}（2-17 の8ビットの\R{階段}{かい|だん}）につなぎます。この\R{部品}{ぶ|ひん}が、\qty{0.01}{V}ごとに1つずつ\R{数}{かず}をふやすとすると、\qty{0.25}{V}は25になります。「\R{温度}{おん|ど}→\R{電圧}{でん|あつ}→\R{数}{かず}」と、2つの\R{箱}{はこ}をつないだ\R{合成}{ごう|せい}\R{関数}{かん|すう}です。

\begin{bkfigure}
\begin{tikzpicture}[x=1mm, y=1mm]
  \bkthermo{(4,-12)}{0.6}{0}{40}{20}{25}
  \node[font=\sffamily\footnotesize, anchor=north] at (4,-23) {\qty{25}{\degreeCelsius}};
  \bkthrdbox{(36,0)}{センサー}{$V=0.01T$}{26}
  \bkthrdbox{(76,0)}{\R{数}{かず}に\R{変}{か}える}{\qty{0.01}{V}で\\1ふえる}{26}
  \bkdisplay{(96,-5.5)}{20}{25}
  \draw[bk link] (10,0) -- (23,0);
  \draw[bk link] (49,0) -- node[bk link label, above] {\qty{0.25}{V}} (63,0);
  \draw[bk link] (89,0) -- (96,0);
\end{tikzpicture}
\end{bkfigure}

\par\addvspace{3pt}\textbf{ブロック\R{線図}{せん|ず}：\R{機械}{き|かい}を「\R{箱}{はこ}と\R{矢印}{や|じるし}」でかく}\par\noindent
\R{機械}{き|かい}を\R{作}{つく}る\R{人}{ひと}は、しくみを「\R{入}{い}れると\R{出}{で}てくる\R{箱}{はこ}」を\R{矢印}{や|じるし}でつないだ\R{図}{ず}でかきます。この\R{図}{ず}を\textbf{ブロック\R{線図}{せん|ず}}といいます。\bkidx{ぶろっくせんず}{ブロック\R{線図}{せん|ず}}たとえば\R{自動車}{じ|どう|しゃ}は、「アクセルをふむ\R{量}{りょう}→エンジンの\R{力}{ちから}→\R{車}{くるま}の\R{速}{はや}さ」と、\R{箱}{はこ}をつないでかけます。\R{第}{だい}9\R{巻}{かん}（9-8）では、この\R{箱}{はこ}の\R{中}{なか}のきまりを\R{式}{しき}で\R{書}{か}き、\R{機械}{き|かい}\R{全体}{ぜん|たい}のふるまいを\R{計算}{けい|さん}します。
\begin{bkfigure}
\begin{tikzpicture}[x=1mm, y=1mm]
  \node[font=\sffamily\footnotesize, align=center] (a) at (6,0) {アクセルを\\ふむ\R{量}{りょう}};
  \node[bk node, minimum width=24mm] (e) at (40,0) {エンジン};
  \node[bk node, minimum width=24mm] (c) at (82,0) {\R{車}{くるま}のからだ};
  \node[font=\sffamily\footnotesize, align=center] (v) at (114,0) {\R{車}{くるま}の\\\R{速}{はや}さ};
  \draw[bk link] (a) -- (e);
  \draw[bk link] (e) -- node[bk link label, above] {\R{力}{ちから}} (c);
  \draw[bk link] (c) -- (v);
\end{tikzpicture}
\end{bkfigure}
\end{butsuri}


%% ---------------- 章末問題 ----------------
\bkendproblems
\bkproblemgroup{\R{基本}{き|ほん}}{1}

\bkproblem $y$ は $x$ の\R{関数}{かん|すう}ですか。\R{関数}{かん|すう}でないものは、$y$ が1つに\R{決}{き}まらない\R{例}{れい}を1つあげましょう。\par
\hspace*{1\zw}（1）\nobreak 1まい80\R{円}{えん}の\R{切手}{きっ|て}を $x$ まい\R{買}{か}ったときの\R{代金}{だい|きん} $y$ \R{円}{えん}\par
\hspace*{1\zw}（2）\nobreak \R{半径}{はん|けい} $x$〔cm〕の\R{円}{えん}の\R{直径}{ちょっ|けい} $y$〔cm〕\par
\hspace*{1\zw}（3）\nobreak まわりの\R{長}{なが}さが $x$〔cm〕の\R{長方形}{ちょう|ほう|けい}の、\R{横}{よこ}の\R{長}{なが}さ $y$〔cm〕
\bkpans{（1）\nobreak \R{関数}{かん|すう}（$y=80x$）\quad（2）\nobreak \R{関数}{かん|すう}（$y=2x$）\quad（3）\nobreak \R{関数}{かん|すう}ではない（\R{例}{れい}：$x=20$ のとき、たて2\,cmなら $y=8$、たて4\,cmなら $y=6$）}

\bkproblem \R{送料}{そう|りょう}が300\R{円}{えん}で、1\R{箱}{はこ}800\R{円}{えん}の\R{紙}{かみ}を $x$ \R{箱}{はこ}\R{注文}{ちゅう|もん}したときの\R{合計}{ごう|けい}を $y$ \R{円}{えん}とします。（1）$y$ を $x$ の\R{式}{しき}で\R{表}{あらわ}しましょう。（2）$x=3$ のときの $y$ を\R{求}{もと}めましょう。（3）\R{合計}{ごう|けい}が4300\R{円}{えん}のとき、\R{箱}{はこ}をいくつ\R{注文}{ちゅう|もん}しましたか。
\bkpans{（1）\nobreak $y=800x+300$\quad（2）\nobreak 2700\R{円}{えん}\quad（3）\nobreak $800x+300=4300$、$800x=4000$、5\R{箱}{はこ}}

\bkproblem $f(x)=-2x+7$ のとき、$f(3)$、$f(-4)$、$f(0.5)$ を\R{求}{もと}めましょう。
\bkpans{$f(3)=1$、$f(-4)=15$、$f(0.5)=6$}

\bkproblem $f(x)=x^2$、$g(x)=4x-3$ のとき、\R{次}{つぎ}の\R{式}{しき}を\R{求}{もと}めましょう。\\（1）$f(-6)$\quad（2）$f(a+2)$\quad（3）$f(5a)$\quad（4）$g(a+h)$
\bkpans{（1）\nobreak 36\quad（2）\nobreak $(a+2)^2$\quad（3）\nobreak $25a^2$\quad（4）\nobreak $4a+4h-3$}

\bkproblem $g(x)=x-2$、$f(u)=u^2$ とします。（1）$f(g(5))$ を\R{求}{もと}めましょう。（2）$y=f(g(x))$ を $x$ の\R{式}{しき}で\R{表}{あらわ}しましょう。（3）$g(f(5))$ を\R{求}{もと}めましょう。
\bkpans{（1）\nobreak $g(5)=3$、$f(3)=9$\quad（2）\nobreak $y=(x-2)^2$\quad（3）\nobreak $f(5)=25$、$g(25)=23$}

\bkproblem $y=-4x$ で、$x$ の\R{変域}{へん|いき}が $-2\leqq x\leqq1$ のとき、$y$ の\R{変域}{へん|いき}を\R{求}{もと}めましょう。
\bkpans{$x=-2$ で $y=8$、$x=1$ で $y=-4$。$-4\leqq y\leqq8$}

\bkproblem \R{次}{つぎ}の\R{比例}{ひ|れい}\R{式}{しき}の $x$ を\R{求}{もと}めましょう。\quad（1）$x:9=4:6$\quad（2）$x:20=3:5$
\bkpans{（1）\nobreak $x=6$\quad（2）\nobreak $x=12$}

\bkproblemgroup{チャレンジ}{2}

\bkproblem \R{摂氏}{せっ|し}を\R{華氏}{か|し}に\R{直}{なお}す\R{関数}{かん|すう}は $f(x)=1.8x+32$ です。（1）$f(-40)$ を\R{求}{もと}めましょう。\R{何}{なに}かふしぎなことに\R{気}{き}づきましたか。（2）\R{華氏}{か|し}で\bkthrdF{212}になるのは、\R{摂氏}{せっ|し}で\R{何}{なん}℃ですか。$f(x)=212$ として\R{方程式}{ほう|てい|しき}を\R{解}{と}きましょう。
\bkpans{（1）\nobreak $f(-40)=1.8\times(-40)+32=-72+32=-40$。\R{摂氏}{せっ|し}と\R{華氏}{か|し}が\R{同}{おな}じ\R{数}{かず}になる\quad（2）\nobreak $1.8x+32=212$、$1.8x=180$、$x=100$ で\qty{100}{\degreeCelsius}}

\bkproblem コインパーキングの\R{料金}{りょう|きん}は「20\R{分}{ぷん}までごとに100\R{円}{えん}」です。とめた\R{時間}{じ|かん}を $x$ \R{分}{ふん}、\R{料金}{りょう|きん}を $y$ \R{円}{えん}とします（$0<x\leqq100$）。（1）50\R{分}{ぷん}とめたときの\R{料金}{りょう|きん}を\R{求}{もと}めましょう。（2）ちょうど80\R{分}{ぷん}とめたときの\R{料金}{りょう|きん}を\R{求}{もと}めましょう。（3）\R{料金}{りょう|きん}が300\R{円}{えん}になるのは、$x$ がどんな\R{範囲}{はん|い}のときですか。
\bkpans{（1）\nobreak 300\R{円}{えん}\quad（2）\nobreak 400\R{円}{えん}\quad（3）\nobreak $40<x\leqq60$}

\bkproblem $g(x)=2x+1$、$f(u)=u^2$ とします。（1）$y=f(g(x))$ と $y=g(f(x))$ を、それぞれ $x$ の\R{式}{しき}で\R{表}{あらわ}しましょう。（2）$x=1$ のとき、それぞれの\R{値}{あたい}を\R{求}{もと}め、ちがうことを\R{確}{たし}かめましょう。
\bkpans{（1）\nobreak $f(g(x))=(2x+1)^2$、$g(f(x))=2x^2+1$\quad（2）\nobreak $f(g(1))=f(3)=9$、$g(f(1))=g(1)=3$}

\bkproblem \R{温度}{おん|ど}センサーは、\R{温度}{おん|ど} $T$〔℃〕を\R{入}{い}れると、\R{電圧}{でん|あつ} $g(T)=0.01T$（V）を\R{出}{だ}します。\R{電圧}{でん|あつ}を\R{数}{かず}に\R{変}{か}える\R{部品}{ぶ|ひん}は、\R{電圧}{でん|あつ} $u$〔V〕を\R{入}{い}れると、\R{数}{かず} $f(u)=100u$ を\R{出}{だ}します。（1）$f(g(25))$ を\R{求}{もと}めましょう。（2）$y=f(g(T))$ を $T$ の\R{式}{しき}で\R{表}{あらわ}しましょう。
\bkpans{（1）\nobreak $g(25)=0.01\times25=0.25$、$f(0.25)=100\times0.25=25$\quad（2）\nobreak $y=100\times0.01T=T$（\R{温度}{おん|ど}と\R{同}{おな}じ\R{数}{かず}が\R{出}{で}てくる）}

\begin{furikaeri}
\begin{checklist}
\item \R{関数}{かん|すう}を\R{見分}{み|わ}け、\R{表}{ひょう}・グラフ・\R{式}{しき}で\R{表}{あらわ}せる。\bkfurref{bk:3-3:kansu}
\item $f(x)$ を\R{読}{よ}み、$f(3)$ や $f(a+h)$ を\R{計算}{けい|さん}できる。\bkfurref{bk:3-3:fx}
\item 2つの\R{関数}{かん|すう}をつないだ\R{合成}{ごう|せい}\R{関数}{かん|すう}を\R{作}{つく}れる。\bkfurref{bk:3-3:gousei}
\end{checklist}
\end{furikaeri}

\bkfukidashi{delta}{\R{関数}{かん|すう}は「\R{入}{い}れると\R{出}{で}てくる\R{箱}{はこ}」。\R{次}{つぎ}の\R{章}{しょう}では、はかった\R{点}{てん}からグラフを\R{自分}{じ|ぶん}でかいて、\R{表}{ひょう}にない\R{値}{あたい}を\R{読}{よ}みとるよ。}
